% !TEX root = ../main.tex
\chapter{Time Series Analysis}\label{ch:timeseries}

\begin{sources}
\small
\begin{tabularx}{\linewidth}{@{}l l L@{}}
\toprule
\textbf{Lecture} & \textbf{Speaker} & \textbf{Content}\\
\midrule
2024 L12 & P.~Kempthorne & Stationary processes; from prices to log returns (S\&P~500, Amazon, crude oil, 10-year yield); sample autocorrelations and the Box--Pierce test; Wold representation theorem; lag operator, impulse response; ARMA, AR($p$) and its stationarity condition, AR(1), Yule--Walker equations, MA($q$); differencing. The lecture ends at slide~22 of 29.\\
2013 L8 (from 26:19) & P.~Kempthorne & The same univariate theory (the first 26 minutes finish regression, \cref{ch:regression}); adds linear-trend reversion versus stochastic trend, ARIMA, maximum likelihood (LIML/FIML) and the information criteria AIC, BIC, HQ.\\
2013 L11 (from 28:35) & P.~Kempthorne & Multivariate time series: cross-covariance matrices, multivariate Wold, VAR($p$) and its VAR(1) form, equation-by-equation least squares and why it is optimal, maximum likelihood, model selection; VAR of US macroeconomic series (Case Study~5). The first 28 minutes are GARCH, \cref{ch:volatility}.\\
2013 L12 & P.~Kempthorne & The fitted macro VAR; cointegration, vector error-correction models, energy futures and crack spreads; linear state-space models (time-varying CAPM, time-varying regression, AR, MA, ARMA); the Kalman filter.\\
\bottomrule
\end{tabularx}

\smallskip
\textbf{Files.} Slides \file{mit18_642_f24_lec11-12.pdf} (29 slides, 73 pages with overlays); 2013 slides
\file{MIT18_S096F13_lecnote8.pdf}, \file{lecnote11.pdf}, \file{lecnote12.pdf}; 2013 Case Study~3
\file{CaseStudy3.pdf} (time series of the US 10-year Treasury yield, in R) and Case Study~5
\file{CaseStudy5.pdf} (VAR models of macroeconomic series, in R).
\textbf{Problem sets.} 2024 PS4 Problems~5--6 (Problems~1--4 are on Brownian motion, \cref{ch:sp2});
2013 PS4 (all six problems); 2013 PS6 Problems~1--2 (Problems~3--5 are portfolio theory,
\cref{ch:portfolio}). The exercises are at the end of the chapter.
\textbf{Not published or not covered.} Slide~5 of the 2024 lecture points to
\file{TimeSeries4plots.pdf} and \file{TimeSeries4acfplots.pdf}, which are not in the download, so the
exploratory plots of \cref{ch09:sec-logret} are described from the spoken commentary. The
crack-spread cointegration note of 2013 L12 exists only in the video. Slides 23--29 of 2024
(trend reversion, ARIMA, maximum likelihood, model selection, Dickey--Fuller test) were not reached in
class; 2013 L8 covered them up to model selection, and the Dickey--Fuller slide was not discussed in
either year.
\end{sources}

\section*{Motivation}
Prices, yields, spreads and macroeconomic indicators are observed \emph{in time}, and the observation
today carries information about tomorrow. Chapter~\ref{ch:sp1} modelled the simplest case, a random walk
with independent steps. Real data have more structure: yield changes are weakly autocorrelated, a credit
spread is pulled back to a long-run level, an interest rate reacts to unemployment and inflation with a
delay, two energy prices drift apart but never far. Time-series analysis supplies the vocabulary:
\begin{itemize}
  \item \emph{stationarity} and the autocorrelation function, to say what ``no structure left'' means;
  \item the \emph{Wold theorem} and the \emph{ARMA} family, to say what structure is possible in
        a stationary series, and \emph{differencing}/unit-root tests to deal with the series that are not
        (\cref{ch09:sec-stat,ch09:sec-wold,ch09:sec-arma,ch09:sec-nonstat});
  \item \emph{VAR} models and \emph{cointegration}, for several series at once
        (\cref{ch09:sec-var,ch09:sec-coint});
  \item \emph{state-space} models and the \emph{Kalman filter}, for hidden quantities such as a time-varying
        beta or a ``true'' price observed with noise (\cref{ch09:sec-ss}).
\end{itemize}
The mathematics is linear algebra and least squares (\cref{ch:linalg,ch:regression}): every model in
this chapter is fitted by regressing the series on its own past, and the stationarity conditions are
statements about the eigenvalues of a matrix. The case studies (10-year Treasury yield and a
three-variable macro VAR) show what the workflow looks like on real data.

% ==================================================================
\section{Stationarity and autocorrelation}\label{ch09:sec-stat}

A \emph{time series} is a stochastic process $\{X_t,\ t\in T\}$ of random variables indexed by the times
of observation, discrete ($T=\{1,2,\dots\}$, the case used almost throughout) or continuous
($T=[0,\infty)$) \ts{24}{12}{0:00}. Its probabilistic behaviour is fixed by the finite-dimensional
distributions $p(x_{t_1},\dots,x_{t_m})$ for every finite set of times \ts{24}{12}{1:09}\ts{13}{8}{27:23}.

\begin{definition}[Strict stationarity]\label{ch09:def-strict}
$\{X_t\}$ is \emph{strictly stationary} if $p(x_{t_1+\tau},\dots,x_{t_m+\tau})=p(x_{t_1},\dots,x_{t_m})$ for all
$\tau$, $m$ and $(t_1,\dots,t_m)$: the joint law is invariant under time translation \ts{24}{12}{2:10}.
\end{definition}

\begin{definition}[Covariance stationarity and ACF]\label{ch09:def-cov}
$\{X_t\}$ is \emph{covariance (weakly, second-order) stationary} if, for all $t$,
\[
  \E X_t=\mu,\qquad \Var X_t=\sigma_X^2<\infty,\qquad \Cov(X_t,X_{t+\tau})=\gamma(\tau)
\]
depend at most on the lag $\tau$ \ts{24}{12}{5:24}\ts{13}{8}{29:32}. The \emph{autocorrelation function}
(ACF) is $\rho(\tau)=\gamma(\tau)/\gamma(0)$ \ts{24}{12}{7:27}.
\end{definition}

Two facts are used constantly. First, $\gamma(-\tau)=\gamma(\tau)$ and $|\rho(\tau)|\le1$. Second, for
every $n$ the Toeplitz matrix $\bigl[\gamma(i-j)\bigr]_{i,j=1}^n$ is the covariance matrix of $(X_1,\dots,X_n)$, hence
positive semi-definite; not every even function $\gamma$ is an autocovariance function.
Strict stationarity together with a finite variance implies covariance stationarity; the converse fails
in general but holds for Gaussian processes, whose law is determined by the first two moments.

\begin{intuition}[Time-translation invariance]
Strict stationarity is the statistician's \emph{equilibrium}: the process looks the same whichever window
one watches, like the noise of a system that has relaxed to a steady state. It is what makes estimation
possible, because a single path can then be averaged over time in place of an ensemble average (the
ergodic hypothesis). Covariance stationarity asks only that the two-point correlation function depend on
time differences, $C(t,t')=\gamma(t-t')$; the Fourier transform of $\gamma$ is the power spectrum, which the
course does not use but which explains why cycles show up as oscillations in $\rho(\tau)$
(\cref{ch09:sec-ar2}). Parameters of a stationary model are constants that can be estimated
\ts{24}{12}{3:12}\ts{13}{8}{28:28}.
\end{intuition}

\subsection{Making financial series stationary}\label{ch09:sec-logret}
Prices are not stationary: they trend, and the size of their fluctuations grows with their level. The
standard remedy is to model the \emph{log return}. Let $Y_t$ be a positive price and
\begin{equation}\label{ch09:eq-logret}
  X_t=\log\frac{Y_t}{Y_{t-1}}=\log(1+R_t),\qquad R_t=\frac{Y_t-Y_{t-1}}{Y_{t-1}},\qquad
  Y_t=Y_0\exp\bigl(X_1+\dots+X_t\bigr),
\end{equation}
so $Y_t$ is an exponential of a random walk with steps $X_t$, and the modelling assumption is that
$\{X_t\}$ is covariance stationary \ts{24}{12}{9:33}\ts{24}{12}{10:33}. For small returns
$X_t\approx R_t-\tfrac12R_t^2$. Taking logs turns exponential growth into linear growth and stabilises
the variance, and differencing removes the linear trend; $X_t$ is then a candidate stationary series.

The lecture illustrated this on four series (the plots themselves are not in the download)
\ts{24}{12}{8:29}:
\begin{itemize}
  \item \emph{S\&P~500}, roughly 2011--2024. The index trends up with sharp COVID drops, so it is not
        stationary in any window; the monthly log returns fluctuate around a level close to zero with
        constant spread. Weekly and daily returns keep a stable mean, but the volatility is clearly not
        constant (extreme in 2020) \ts{24}{12}{11:42}\ts{24}{12}{13:48}. A fitted normal density
        underestimates both the mass near the mean and the tails: the returns are
        \emph{leptokurtic} (``slender''), as in \cref{ch03:sec-heavy}. Volatility clustering is the subject of
        \cref{ch:volatility}.
  \item \emph{Amazon}. The price shows strong exponential growth; on the log scale the variability is more
        homogeneous, and the monthly and weekly log returns look stationary with a better normal fit than the
        index. Daily returns show a few spikes that no Gaussian model would produce
        \ts{24}{12}{15:56}\ts{24}{12}{16:56}.
  \item \emph{Crude oil futures}. In April 2020 the front-month contract went \emph{negative}; a holder of
        100 contracts woke up owing the broker 100 times a negative price, and some brokers could not
        even report positions because their systems did not allow negative prices. The logarithm is
        undefined there, so the lecture's plots exclude those days \ts{24}{12}{18:00}\ts{24}{12}{19:03}.
        Log returns require positive prices (spreads and rates can be modelled directly).
  \item \emph{10-year US Treasury yield}, from above $3\%$ to below $1\%$ after COVID and rising again.
        Strictly, ``log return'' is the wrong name for a percentage change of a yield, but the same
        transformation gives unimodal, roughly symmetric increments, and the normal model fits better than
        for equity returns \ts{24}{12}{20:05}\ts{24}{12}{21:12}\ts{24}{12}{22:16}.
\end{itemize}
Modelling $Y_t$ as an exponential random walk means that $S_t=X_1+\dots+X_t$ is a sum of steps. In the
simplest model the steps are i.i.d.\ (\cref{ch04:sec-rw}); the models of this chapter let them depend
on their own past \ts{24}{12}{23:18}.

\subsection{Sample autocorrelation and tests of whiteness}\label{ch09:sec-sacf}
Given observations $x_1,\dots,x_T$ the \emph{sample autocorrelation at lag $k$} is
\begin{equation}\label{ch09:eq-sacf}
  \hat r_k=\frac{\sum_{t=1}^{T-k}(x_t-\bar x)(x_{t+k}-\bar x)}{\sum_{t=1}^{T}(x_t-\bar x)^2}
\end{equation}
\ts{24}{12}{24:25}, an estimate of $\rho(k)$. If the series is white noise, then $\hat r_k$ is
approximately $N(0,1/T)$ (Bartlett); the lecture writes the variance as $1/(T-k)$, which is the same for
$T\gg k$ \ts{24}{12}{25:35}. Hence the test of $H_0:\rho(k)=0$ rejects at the $5\%$ level when
\begin{equation}\label{ch09:eq-band}
  |\hat r_k|>1.96\sqrt{\tfrac1{T-k}},
\end{equation}
which is the pair of dashed lines that R's \texttt{acf} draws around zero \ts{24}{12}{26:42}. For
a daily series with $T=500$ the bands are $\pm0.088$; a single autocorrelation of $0.09$ is
already significant. The lecture saw a large negative monthly autocorrelation at lag~1 for the
S\&P~500 and, on weekly data, significant negative correlations around lag~5, which
a student read as \emph{mean reversion} (price changes that are too high tend to be reversed)
\ts{24}{12}{27:42}. Amazon showed no significant autocorrelation at any frequency, crude oil a negative
lag-2 and positive lag-3 weekly pattern, and Treasury yields strong low-lag dependence
\ts{24}{12}{29:50}\ts{24}{12}{30:58}.

\paragraph{Testing many lags at once.} If $\rho(1)=\dots=\rho(K)=0$ the $\sqrt T\hat r_j$ are
asymptotically independent $N(0,1)$, so
\begin{equation}\label{ch09:eq-bp}
  \mathrm{BP}=T\sum_{j=1}^{K}\hat r_j^{\,2}\ \approx\ \chi^2_K ,
\end{equation}
the \emph{Box--Pierce} statistic; reject the joint null when $\mathrm{BP}$ exceeds the $\chi^2_K$
quantile ($18.31$ for $K=10$ at $5\%$) \ts{24}{12}{36:10}\ts{24}{12}{37:18}. Two refinements, not
discussed in class: the Ljung--Box variant $T(T+2)\sum_j\hat r_j^2/(T-j)$ has a better small-sample
$\chi^2$ approximation; and when the test is applied to the residuals of a fitted ARMA$(p,q)$ model the
degrees of freedom drop to $K-p-q$.

\begin{finance}[Residual whiteness as a model check]
A time-series model writes $Y_t=\hat Y_t+e_t$. It is \emph{adequate} when the residuals have mean zero,
constant variance and no autocorrelation, i.e.\ they are white noise: nothing predictable is left
\ts{24}{12}{31:59}\ts{24}{12}{33:03}. If in addition they are Gaussian, the model delivers not only point
forecasts but valid \emph{prediction intervals} \ts{24}{12}{34:03}. The ACF of the residuals, and the
Box--Pierce test on it, are therefore the standard diagnostic: a model whose residuals still show
significant autocorrelation ``is not done yet'' \ts{24}{12}{38:24}. The flip side, developed in
\cref{ch:volatility}, is that uncorrelated does not mean independent: returns can be white noise while
their squares are strongly autocorrelated.
\end{finance}

% ==================================================================
\section{The Wold representation theorem}\label{ch09:sec-wold}

\begin{theorem}[Wold representation]\label{ch09:thm-wold}
Every zero-mean covariance-stationary series $\{X_t\}$ can be written
\begin{equation}\label{ch09:eq-wold}
  X_t=V_t+S_t,\qquad S_t=\sum_{i=0}^{\infty}\psi_i\,\eta_{t-i},
\end{equation}
where $\psi_0=1$, $\sum_i\psi_i^2<\infty$; $\{\eta_t\}$ is \emph{white noise}, i.e.\ $\E\eta_t=0$,
$\E\eta_t^2=\sigma^2$, $\E\eta_t\eta_s=0$ for $t\ne s$; $\{V_t\}$ is \emph{linearly deterministic} (a linear
combination of its own past with constant coefficients predicts it exactly); and
$\E(\eta_tV_s)=0$ for all $t,s$ \ts{24}{12}{38:24}\ts{13}{8}{31:40}.
\end{theorem}

The theorem says that once the deterministic part is removed, \emph{every} stationary series is a moving
average of white noise, possibly of infinite order. The white noise $\eta_t=X_t-P(X_t\mid X_{t-1},X_{t-2},\dots)$
is the \emph{innovation}: the part of $X_t$ that cannot be predicted linearly from the past, the ``news'' arriving
at $t$ \ts{13}{11}{43:32}\ts{13}{11}{44:38}. An example of a linearly deterministic process is a
sinusoid, $V_t=A_1\cos\omega t+A_2\sin\omega t$ with random $A_1,A_2$: it obeys the exact recursion
$V_t=2\cos\omega\,V_{t-1}-V_{t-2}$ (from $\cos\omega t=2\cos\omega\cos\omega(t-1)-\cos\omega(t-2)$), so two
observations fix it for ever; sums of sinusoids of different frequencies (a Fourier series) are also linearly
predictable \ts{24}{12}{40:30}\ts{24}{12}{41:32}\ts{24}{12}{42:32}.

\begin{proof}[Sketch of proof]
Work in the Hilbert space $L^2$ of square-integrable random variables with inner product $\langle U,W\rangle=\E UW$.
Let $H_t$ be the closed linear span of $\{X_s:s\le t\}$ and $P_{H}$ the orthogonal projection on $H$. Define
$\eta_t=X_t-P_{H_{t-1}}X_t$. Stationarity makes the time shift a unitary map, so $\sigma^2=\E\eta_t^2$ does not depend on
$t$ (we assume $\sigma^2>0$; otherwise $X$ is itself deterministic). Since $\eta_s\in H_s\subset H_{t-1}$ for $s<t$
and $\eta_t\perp H_{t-1}$, the $\eta_t$ are uncorrelated. Put $\psi_j=\E[X_t\eta_{t-j}]/\sigma^2$, independent of $t$; then
$\psi_0=\E[X_t\eta_t]/\sigma^2=\E[\eta_t^2]/\sigma^2=1$, because $X_t-\eta_t\in H_{t-1}$. The projection of $X_t$ on
$\mathrm{span}\{\eta_t,\dots,\eta_{t-J+1}\}$ is $\sum_{j<J}\psi_j\eta_{t-j}$, and Bessel's inequality gives
$\sigma^2\sum_j\psi_j^2\le\Var X_t<\infty$. Hence $S_t=\sum_{j\ge0}\psi_j\eta_{t-j}$ converges in $L^2$, and
$V_t=X_t-S_t$ is orthogonal to every $\eta_s$ with $s\le t$. One checks that $V_t\in\bigcap_sH_s$, the
\emph{remote past} of the process, so $V_t$ is determined exactly, by a linear combination, from
arbitrarily remote values: it is linearly deterministic.
\end{proof}

\begin{intuition}[What the theorem does and does not say]
The proof shows why the result holds: $\eta_t$ is the component of $X_t$ orthogonal to the past, and
expanding $X_t$ in the orthogonal family $\eta_t,\eta_{t-1},\dots$ is a Gram--Schmidt procedure. The
innovations are only \emph{uncorrelated}, not independent: the theorem describes the best \emph{linear}
prediction, and a process such as a GARCH model (\cref{ch:volatility}) is covered by it although its
innovations are dependent through their variances. Covariance stationarity of the sum
$S_t$ needs only $\sum\psi_i^2<\infty$; that is why the condition appears \ts{24}{12}{50:07}.
\end{intuition}

\subsection{Wold's theorem as a modelling recipe}\label{ch09:sec-wold-recipe}
The theorem suggests how to specify a stationary model for data $x_0,\dots,x_T$
\ts{24}{12}{44:41}\ts{13}{8}{34:48}. Fix a number $p$ of lags to represent the linearly deterministic part, and take an
estimation window of size $n$ (after de-trending so that the series has a constant mean level). Set
\[
  \bm y=\begin{pmatrix}y_1\\ \vdots\\ y_n\end{pmatrix},\qquad
  Z=\begin{pmatrix}1&y_0&y_{-1}&\cdots&y_{-(p-1)}\\ 1&y_1&y_0&\cdots&y_{-(p-2)}\\
  \vdots&\vdots&&&\vdots\\ 1&y_{n-1}&y_{n-2}&\cdots&y_{n-p}\end{pmatrix},
\]
so the first $p$ rows reach back before the start of the response. The linear projection on $p$ lags is the least-squares fit
\[
  \hat{\bm y}^{(p)}=Z(Z\T Z)^{-1}Z\T\bm y=\hat P(Y_t\mid Y_{t-1},\dots,Y_{t-p}),\qquad
  \hat{\bm\varepsilon}^{(p)}=\bm y-\hat{\bm y}^{(p)},
\]
the hat matrix of \cref{ch05:prop-hat}. The residual series is then modelled as a moving average
$\varepsilon_t^{(p)}=\sum_i\psi_i\eta_{t-i}$, giving estimates of the $\psi_i$ and of the innovations $\eta_t$
\ts{24}{12}{45:41}\ts{24}{12}{46:45}. Two diagnostics decide whether to stop:
\begin{itemize}
  \item the residual should be orthogonal to lags of $Y$ beyond $p$; if it is not, increase $p$;
  \item the estimated innovations should look like white noise (\cref{ch09:sec-sacf}); if not, change the
        specification of the moving average or add deterministic regressors \ts{24}{12}{47:50}.
\end{itemize}
As $p\to\infty$ the projection converges to $P(Y_t\mid Y_{t-1},Y_{t-2},\dots)$, but a sample of size $n$
supports only $p/n\to0$, so useful models have a modest $p$ or a moving average with a
parsimonious parametrisation \ts{24}{12}{48:55}\ts{13}{8}{38:55}. This is exactly why ARMA models
exist (\cref{ch09:sec-arma}).

\subsection{Lag operator, impulse response, invertibility}\label{ch09:sec-lag}
The \emph{lag} (backshift) operator $L$ shifts a series back one period: $LX_t=X_{t-1}$, $L^0=\mathrm{id}$,
$L^nX_t=X_{t-n}$; its inverse $L^{-n}X_t=X_{t+n}$ shifts forward \ts{24}{12}{52:09}\ts{13}{8}{40:01}. Polynomials in $L$ multiply
and factorise like ordinary polynomials in a complex variable $z$, and the Wold representation reads
\begin{equation}\label{ch09:eq-woldL}
  X_t=\psi(L)\eta_t+V_t,\qquad \psi(L)=\sum_{i\ge0}\psi_iL^i .
\end{equation}
The \emph{impulse response function} is $\mathrm{IR}(j)=\partial X_t/\partial\eta_{t-j}=\psi_j$, the effect of a
shock $j$ periods ago on the process today \ts{24}{12}{54:20}; graphed against $j$ it displays delayed
responses. If the Federal Reserve chairman changes interest rates and the variable of interest is GNP,
$\psi_j$ is the effect $j$ quarters later \ts{13}{8}{43:11}\ts{13}{8}{44:11}. The \emph{long-run cumulative response} is
$\sum_j\psi_j=\psi(1)$, the level to which the process eventually moves after one unit innovation
\ts{24}{12}{56:28}\ts{13}{8}{45:11}.

If $\psi(L)$ is \emph{invertible}, i.e.\ there is $\psi^{-1}(L)=\sum_i\psi^*_iL^i$ with
$\psi^{-1}(L)\psi(L)=L^0$, then for $V_t=0$
\begin{equation}\label{ch09:eq-arinf}
  \psi^{-1}(L)X_t=\eta_t\quad\Longleftrightarrow\quad X_t=\sum_{i\ge1}\pi_iX_{t-i}+\eta_t,\qquad \pi_i=-\psi^*_i,
\end{equation}
an infinite-order autoregression \ts{24}{12}{57:31}\ts{24}{12}{58:35}\ts{13}{8}{46:17}. The coefficients of the
least-squares projection on $p$ lags in \cref{ch09:sec-wold-recipe} converge to these $\pi_i$
\ts{13}{8}{47:20}. The basic tool for inverting is the geometric series: for a complex number $\lambda$ with
$|\lambda|>1$,
\begin{equation}\label{ch09:eq-geom}
  \bigl(1-\lambda^{-1}L\bigr)^{-1}=\sum_{i\ge0}\lambda^{-i}L^i,
\end{equation}
which converges (on bounded stationary series) precisely because the coefficients $\lambda^{-i}$ decay
\ts{24}{12}{1:06:02}.

% ==================================================================
\section{ARMA models}\label{ch09:sec-arma}

\begin{definition}[ARMA$(p,q)$]\label{ch09:def-arma}
$\{X_t\}$ follows an \emph{autoregressive moving-average} model of orders $(p,q)$ if
\begin{equation}\label{ch09:eq-arma}
  X_t-\mu=\phi_1(X_{t-1}-\mu)+\dots+\phi_p(X_{t-p}-\mu)+\eta_t+\theta_1\eta_{t-1}+\dots+\theta_q\eta_{t-q},
\end{equation}
where $\{\eta_t\}$ is white noise $\mathrm{WN}(0,\sigma^2)$. With the polynomials
$\phi(L)=1-\phi_1L-\dots-\phi_pL^p$ and $\theta(L)=1+\theta_1L+\dots+\theta_qL^q$ this reads
\[
  \phi(L)(X_t-\mu)=\theta(L)\eta_t,\qquad\text{and}\qquad X_t=\mu+\psi(L)\eta_t,\quad \psi(L)=\frac{\theta(L)}{\phi(L)}
\]
is the Wold decomposition, when $\phi(L)$ is invertible \ts{24}{12}{59:35}\ts{13}{8}{46:17}\ts{13}{8}{48:22}.
\end{definition}

The sign convention of the moving-average part is ``$+\theta_j$'' (as in R); some texts use ``$-\theta_j$''.
ARMA models are \emph{parsimonious}: $p+q$ parameters generate an infinite Wold expansion
\ts{24}{12}{1:00:39}. We first look at the two pure cases.

\subsection{Autoregressive models AR($p$)}\label{ch09:sec-ar}
Set $q=0$: $\phi(L)(X_t-\mu)=\eta_t$, that is,
\begin{equation}\label{ch09:eq-arp}
  X_t=c+\sum_{j=1}^p\phi_jX_{t-j}+\eta_t,\qquad c=\mu\,\phi(1)=\mu\Bigl(1-\sum_j\phi_j\Bigr),
\end{equation}
a \emph{linear regression of the series on its own lags}, with white-noise errors
\ts{24}{12}{1:01:47}\ts{24}{12}{1:02:49}\ts{13}{8}{50:34}. (Taking expectations of \eqref{ch09:eq-arp} gives $\mu=c+\mu\sum\phi_j$,
whence $c=\mu\phi(1)$.)

\begin{theorem}[Stationarity of AR($p$)]\label{ch09:thm-ar-stat}
The equation $\phi(L)(X_t-\mu)=\eta_t$ has a covariance-stationary solution of the form
$X_t-\mu=\sum_{i\ge0}\psi_i\eta_{t-i}$ if and only if every root $\lambda_j$ of the
\emph{characteristic equation} $\phi(z)=1-\phi_1z-\dots-\phi_pz^p=0$ lies outside the closed unit
disc, $|\lambda_j|>1$, $j=1,\dots,p$ \ts{24}{12}{1:03:51}\ts{13}{8}{52:35}.
\end{theorem}
\begin{proof}
Factorise $\phi(z)=\prod_{j=1}^p(1-z/\lambda_j)$ (the constant term is~1). If all $|\lambda_j|>1$,
each factor is inverted by the geometric series \eqref{ch09:eq-geom}, whose coefficients $\lambda_j^{-i}$ decay
geometrically, and the product $\phi^{-1}(L)=\prod_j(1-L/\lambda_j)^{-1}$ of absolutely summable series is
absolutely summable. Thus $\psi(L)=\phi^{-1}(L)$ has $\sum|\psi_i|<\infty$ and $X_t=\mu+\psi(L)\eta_t$ is
covariance stationary \ts{24}{12}{1:05:00}\ts{24}{12}{1:07:09}. Conversely, a root with $|\lambda|<1$ makes
the coefficients $\lambda^{-i}$ grow, so the expansion diverges (the process is explosive), and a root on
the circle $|\lambda|=1$ gives a unit root (the random walk is the case $p=1$, $\lambda=1$), whose variance
grows with $t$; in neither case is there a stationary solution driven by the past.
\end{proof}

\begin{remark}[Roots and eigenvalues]\label{ch09:rem-companion}
The condition is often stated on the reciprocals $1/\lambda_j$, the roots of $z^p-\phi_1z^{p-1}-\dots-\phi_p=0$,
which are the eigenvalues of the \emph{companion matrix}
\[
  F=\begin{pmatrix}\phi_1&\phi_2&\cdots&\phi_{p-1}&\phi_p\\ 1&0&\cdots&0&0\\ 0&1&&0&0\\ \vdots&&\ddots&&\vdots\\ 0&0&\cdots&1&0\end{pmatrix};
\]
stationarity is $|1/\lambda_j|<1$ for all $j$. Software conventions differ: R's \texttt{vars} package reports
the moduli of the eigenvalues (stationary if all $<1$), whereas \texttt{polyroot} returns roots of $\phi(z)$
(stationary if all $>1$) \ts{13}{12}{1:08}. The proof of the correspondence is given for the vector case in
\cref{ch09:prop-companion}.
\end{remark}

\begin{figure}[ht]
\centering
\begin{tikzpicture}[scale=1.15,>=Stealth]
  \draw[->] (-2.4,0)--(2.4,0) node[right] {$\mathrm{Re}\,z$};
  \draw[->] (0,-2.1)--(0,2.1) node[above] {$\mathrm{Im}\,z$};
  \draw[thick,cscol,fill=cscol!8] (0,0) circle (1);
  \node[cscol] at (0.05,0.5) {\footnotesize\begin{tabular}{c}roots here:\\ explosive\end{tabular}};
  \foreach \x/\y in {1.7/0.6,1.7/-0.6,-1.5/0,0/1.6}{\node[intcol] at (\x,\y) {\large$\times$};}
  \node[intcol,align=left,font=\footnotesize] at (1.55,1.55) {stationary\\[-2pt](all roots\\[-2pt]outside)};
  \node[thmcol] at (0.7,-0.45) {\large$\times$};
  \node[thmcol,font=\footnotesize] at (-1.05,-1.55) {one root inside};
  \draw[thmcol,->] (-1.0,-1.35)--(0.62,-0.5);
  \node[font=\footnotesize] at (1.2,-0.2) {$1$};
\end{tikzpicture}
\caption{Roots $\lambda_j$ of $\phi(z)=0$ in the complex plane. Green crosses: a stationary AR model, all roots
outside the unit circle. Red cross: a single root inside makes the process explosive.}\label{ch09:fig-roots}
\end{figure}

\subsubsection{AR(1)}\label{ch09:sec-ar1}
For $X_t-\mu=\phi(X_{t-1}-\mu)+\eta_t$ the characteristic equation is $1-\phi z=0$ with root $\lambda=1/\phi$:
the process is covariance stationary iff $|\phi|<1$ \ts{24}{12}{1:08:12}. Iterating,
$X_t-\mu=\sum_{j\ge0}\phi^j\eta_{t-j}$ (the Wold decomposition, $\psi_j=\phi^j$), and therefore
\begin{equation}\label{ch09:eq-ar1}
  \E X_t=\mu,\qquad \gamma(0)=\sigma_X^2=\frac{\sigma^2}{1-\phi^2},\qquad
  \gamma(j)=\phi^{\,j}\,\sigma_X^2,\qquad \rho(j)=\phi^{\,j}
\end{equation}
\ts{24}{12}{1:08:12}\ts{13}{8}{56:46}\ts{13}{8}{58:54}. (Multiply $X_t-\mu=\phi(X_{t-1}-\mu)+\eta_t$ by $X_{t-j}-\mu$
and take expectations, or sum the geometric series.) The variance exceeds the innovation variance $\sigma^2$ for every
$\phi\neq0$, and depends on $\phi$ only through $\phi^2$: positive and negative autocorrelation inflate it equally
\ts{13}{8}{57:50}. The regimes:
\begin{itemize}
  \item $0<\phi<1$: \emph{exponential mean reversion} to $\mu$, all correlations positive;
  \item $-1<\phi<0$: \emph{oscillating} exponential mean reversion, correlations alternate in sign
        \ts{24}{12}{1:09:15};
  \item $\phi=1$: the Wold decomposition does not exist, $X_t$ is a random walk, not stationary;
  \item $|\phi|>1$: explosive.
\end{itemize}

\begin{figure}[ht]
\centering
\begin{tikzpicture}
\begin{axis}[width=0.47\linewidth,height=4.6cm,xlabel={lag $j$},ylabel={$\rho(j)$},ymin=-1.05,ymax=1.05,
  xmin=-0.5,xmax=12.5,ytick={-1,-0.5,0,0.5,1},title={AR(1), $\phi=0.8$},title style={font=\footnotesize},grid=major,grid style={black!10}]
  \addplot+[ycomb,thick,mark=*,mark size=1.5pt,color=defcol] coordinates {(0,1.0) (1,0.8) (2,0.64) (3,0.512) (4,0.4096) (5,0.3277) (6,0.2621) (7,0.2097) (8,0.1678) (9,0.1342) (10,0.1074) (11,0.0859) (12,0.0687)};
\end{axis}
\end{tikzpicture}\hfill
\begin{tikzpicture}
\begin{axis}[width=0.47\linewidth,height=4.6cm,xlabel={lag $j$},ymin=-1.05,ymax=1.05,
  xmin=-0.5,xmax=12.5,ytick={-1,-0.5,0,0.5,1},title={AR(1), $\phi=-0.8$},title style={font=\footnotesize},grid=major,grid style={black!10}]
  \addplot+[ycomb,thick,mark=*,mark size=1.5pt,color=thmcol] coordinates {(0,1.0) (1,-0.8) (2,0.64) (3,-0.512) (4,0.4096) (5,-0.3277) (6,0.2621) (7,-0.2097) (8,0.1678) (9,-0.1342) (10,0.1074) (11,-0.0859) (12,0.0687)};
\end{axis}
\end{tikzpicture}
\caption{Theoretical autocorrelation $\rho(j)=\phi^j$ of an AR(1): geometric decay for $\phi>0$, geometric
decay with alternating sign for $\phi<0$.}\label{ch09:fig-acf-ar1}
\end{figure}

\begin{figure}[ht]
\centering
\begin{tikzpicture}
\begin{axis}[width=0.95\linewidth,height=5.2cm,xlabel={$t$},xmin=0,xmax=149,legend pos=south west,legend style={font=\footnotesize},
  grid=major,grid style={black!10}]
  \addplot[thick,color=defcol] coordinates {(0,0.0) (1,-2.556) (2,-1.754) (3,-2.059) (4,-2.203) (5,-2.088) (6,-3.795) (7,-3.457) (8,-3.804) (9,0.09) (10,0.302) (11,-0.096) (12,-0.363) (13,-0.976) (14,-1.885) (15,-1.993) (16,-1.212) (17,-1.269) (18,-0.121) (19,-0.303) (20,-0.233) (21,1.348) (22,1.691) (23,0.932) (24,0.609) (25,1.058) (26,2.835) (27,2.14) (28,1.575) (29,2.341) (30,1.104) (31,0.646) (32,1.432) (33,1.798) (34,1.619) (35,2.047) (36,-1.089) (37,0.096) (38,-0.878) (39,-2.415) (40,-1.776) (41,-0.809) (42,-1.133) (43,-2.039) (44,-1.707) (45,-1.504) (46,0.127) (47,0.856) (48,0.921) (49,1.895) (50,1.405) (51,0.268) (52,0.812) (53,1.273) (54,0.867) (55,-0.046) (56,0.19) (57,-2.332) (58,-1.292) (59,-0.607) (60,-2.155) (61,-1.77) (62,-2.469) (63,-1.341) (64,-3.174) (65,-3.613) (66,-2.361) (67,-0.851) (68,-2.881) (69,-2.947) (70,-2.177) (71,-2.46) (72,-0.5) (73,-1.616) (74,-1.019) (75,-1.915) (76,-0.222) (77,-0.21) (78,-0.551) (79,-2.186) (80,-0.177) (81,0.603) (82,1.266) (83,2.214) (84,2.231) (85,1.257) (86,0.268) (87,-0.572) (88,0.884) (89,-0.709) (90,-1.199) (91,-1.341) (92,-0.915) (93,-0.202) (94,-1.421) (95,-2.938) (96,-2.502) (97,-0.913) (98,-0.019) (99,0.2) (100,-0.147) (101,0.168) (102,-0.101) (103,0.732) (104,-0.173) (105,-0.012) (106,-0.121) (107,0.44) (108,0.599) (109,3.059) (110,4.099) (111,4.981) (112,2.194) (113,1.525) (114,0.687) (115,1.117) (116,-1.33) (117,0.044) (118,1.105) (119,-0.363) (120,-1.287) (121,-1.895) (122,-1.568) (123,-0.692) (124,1.46) (125,1.044) (126,1.655) (127,1.562) (128,3.087) (129,3.366) (130,4.23) (131,2.518) (132,1.948) (133,0.842) (134,2.221) (135,2.545) (136,1.858) (137,1.127) (138,1.443) (139,0.525) (140,-0.485) (141,0.069) (142,2.522) (143,1.898) (144,1.057) (145,-0.273) (146,-1.567) (147,-0.807) (148,0.165) (149,0.15)};
  \addplot[thick,color=thmcol] coordinates {(0,0.0) (1,-2.556) (2,-2.138) (3,-2.705) (4,-3.158) (5,-3.374) (6,-5.394) (7,-5.626) (8,-6.491) (9,-3.168) (10,-2.942) (11,-3.295) (12,-3.576) (13,-4.244) (14,-5.299) (15,-5.69) (16,-5.208) (17,-5.446) (18,-4.489) (19,-4.688) (20,-4.664) (21,-3.118) (22,-2.573) (23,-3.079) (24,-3.261) (25,-2.721) (26,-0.786) (27,-1.055) (28,-1.299) (29,-0.297) (30,-1.183) (31,-1.475) (32,-0.592) (33,-0.012) (34,0.08) (35,0.75) (36,-2.078) (37,-1.057) (38,-2.017) (39,-3.685) (40,-3.409) (41,-2.708) (42,-3.153) (43,-4.23) (44,-4.203) (45,-4.256) (46,-2.851) (47,-2.103) (48,-1.909) (49,-0.798) (50,-1.003) (51,-1.929) (52,-1.345) (53,-0.763) (54,-0.977) (55,-1.76) (56,-1.531) (57,-4.025) (58,-3.335) (59,-2.843) (60,-4.482) (61,-4.421) (62,-5.385) (63,-4.628) (64,-6.662) (65,-7.577) (66,-6.867) (67,-5.711) (68,-7.869) (69,-8.367) (70,-8.039) (71,-8.648) (72,-7.057) (73,-8.248) (74,-7.894) (75,-8.942) (76,-7.536) (77,-7.558) (78,-7.93) (79,-9.648) (80,-7.967) (81,-7.214) (82,-6.46) (83,-5.322) (84,-4.973) (85,-5.612) (86,-6.413) (87,-7.213) (88,-5.843) (89,-7.303) (90,-7.899) (91,-8.221) (92,-7.996) (93,-7.421) (94,-8.67) (95,-10.4) (96,-10.404) (97,-9.191) (98,-8.434) (99,-8.218) (100,-8.535) (101,-8.242) (102,-8.485) (103,-7.668) (104,-8.462) (105,-8.328) (106,-8.439) (107,-7.896) (108,-7.671) (109,-5.121) (110,-3.622) (111,-2.126) (112,-4.165) (113,-4.505) (114,-5.114) (115,-4.581) (116,-6.86) (117,-5.686) (118,-4.619) (119,-5.921) (120,-6.899) (121,-7.701) (122,-7.657) (123,-7.016) (124,-4.968) (125,-5.166) (126,-4.398) (127,-4.243) (128,-2.483) (129,-1.741) (130,-0.372) (131,-1.45) (132,-1.642) (133,-2.456) (134,-0.951) (135,-0.293) (136,-0.599) (137,-1.051) (138,-0.566) (139,-1.268) (140,-2.199) (141,-1.717) (142,0.746) (143,0.5) (144,-0.056) (145,-1.227) (146,-2.562) (147,-2.037) (148,-1.186) (149,-1.177)};
  \legend{{AR(1), $\phi=0.85$},{random walk ($\phi=1$)}}
\end{axis}
\end{tikzpicture}
\caption{Two paths driven by the same simulated white noise (seed~3). The AR(1) with $\phi=0.85$ keeps
returning to $\mu=0$; the random walk wanders (its variance grows like $t\sigma^2$).}\label{ch09:fig-ar1-rw}
\end{figure}

\begin{finance}[AR(1) in practice: mean reversion and Vasicek]\label{ch09:fin-vasicek}
The AR(1) is the discrete-time version of the Ornstein--Uhlenbeck (Vasicek) process
$\mathrm dX=-\kappa(X-\mu)\,\mathrm dt+\sigma_0\,\mathrm dW$ (\cref{ch:sde}): sampled every $h$, the exact
transition is an AR(1) with $\phi=e^{-\kappa h}$ and noise variance $\sigma_0^2(1-e^{-2\kappa h})/(2\kappa)$, consistent with
$\gamma(0)=\sigma_0^2/(2\kappa)$ in \eqref{ch09:eq-ar1}. The half-life of a shock is
$h\ln2/(-\ln\phi)$. Candidates for such a model are variables that should not wander too far: interest
rates, interest-rate spreads, real exchange rates and valuation ratios (dividend-to-price, earnings-to-price)
\ts{24}{12}{1:10:15}\ts{24}{12}{1:11:16}\ts{13}{8}{58:54}. The caveat of the 2013 lecture: this is usually true over short
periods; over long ones the level $\mu$ to which the series reverts drifts, so the model must be re-estimated
\ts{13}{8}{59:59}.
\end{finance}

\subsubsection{Yule--Walker equations}\label{ch09:sec-yw}
Multiply the AR($p$) equation $X_t-\mu=\sum_i\phi_i(X_{t-i}-\mu)+\eta_t$ by $X_{t-j}-\mu$ and take expectations. Since
$\E[\eta_t(X_{t-j}-\mu)]=0$ for $j\ge1$ and $=\sigma^2$ for $j=0$ (the noise is uncorrelated with the past),
\begin{equation}\label{ch09:eq-yw}
  \gamma(j)=\sum_{i=1}^p\phi_i\,\gamma(j-i)+\sigma^2\,\delta_{j0},\qquad j=0,1,2,\dots
\end{equation}
\ts{24}{12}{1:11:16}\ts{24}{12}{1:12:16}\ts{13}{8}{1:02:01} (the slide misprints the second lag as $X_{t-1}$; the
lecturer corrected it aloud). Dividing by $\gamma(0)$, $\rho(j)=\sum_i\phi_i\rho(j-i)$ for $j\ge1$. The equations
$j=1,\dots,p$ form the \emph{Yule--Walker system}
\begin{equation}\label{ch09:eq-ywsys}
  \begin{pmatrix}\gamma(1)\\ \gamma(2)\\ \vdots\\ \gamma(p)\end{pmatrix}=
  \begin{pmatrix}\gamma(0)&\gamma(1)&\cdots&\gamma(p-1)\\ \gamma(1)&\gamma(0)&\cdots&\gamma(p-2)\\
  \vdots&&\ddots&\vdots\\ \gamma(p-1)&\gamma(p-2)&\cdots&\gamma(0)\end{pmatrix}
  \begin{pmatrix}\phi_1\\ \phi_2\\ \vdots\\ \phi_p\end{pmatrix},
\end{equation}
using $\gamma(-j)=\gamma(j)$. The matrix is Toeplitz (constant along diagonals), a question the 2013 lecturer put to the
class \ts{13}{8}{1:03:02}. The zeroth equation then gives the noise variance
\begin{equation}\label{ch09:eq-ywsigma}
  \sigma^2=\gamma(0)-\sum_{j=1}^p\phi_j\gamma(j).
\end{equation}
\textbf{Estimation.} Replace the $\gamma(j)$ by sample autocovariances $\hat\gamma(j)$ and solve
\eqref{ch09:eq-ywsys}: these are the \emph{Yule--Walker estimates} \ts{24}{12}{1:13:20}. They are an application of the
\emph{method of moments}: equate sample and theoretical moments, and solve for the parameters. This is the tool of choice when the likelihood is hard to write
down or to optimise; econometrics is rich in such applications \ts{24}{12}{1:14:24}\ts{13}{8}{1:05:09}. (Two remarks beyond
the course: the system is solved efficiently by the Levinson--Durbin recursion, and with the divisor $T$ in $\hat\gamma$ the
fitted AR model is always stationary.)

\begin{example}[Yule--Walker with two lags]\label{ch09:ex-yw}
Suppose the sample autocorrelations of a series are $\hat\rho(1)=0.5$ and $\hat\rho(2)=0.2$. The $p=2$ system
$\bigl(\begin{smallmatrix}1&0.5\\0.5&1\end{smallmatrix}\bigr)\bigl(\begin{smallmatrix}\phi_1\\\phi_2\end{smallmatrix}\bigr)
=\bigl(\begin{smallmatrix}0.5\\0.2\end{smallmatrix}\bigr)$ has determinant $0.75$ and solution
$\phi_1=(0.5-0.5\cdot0.2)/0.75=8/15\approx0.533$ and $\phi_2=(0.2-0.25)/0.75=-1/15\approx-0.067$. The fitted process is
stationary: $\phi_1+\phi_2=0.467<1$, $\phi_2-\phi_1=-0.6<1$, $|\phi_2|<1$ (\cref{ch09:sec-ar2}). On a simulated AR(2)
with $\phi=(1,-0.5)$ and $n=2000$ the Yule--Walker estimates were $(1.017,-0.515)$ with $\hat\sigma^2=0.960$
(true value~1).
\end{example}

\subsubsection{Partial autocorrelation}\label{ch09:sec-pacf}
The \emph{partial autocorrelation} at lag $k$, $\alpha(k)$, is the correlation between $X_t$ and $X_{t-k}$ \emph{after
removing the linear effect of the intermediate values} $X_{t-1},\dots,X_{t-k+1}$
\ts{24}{12}{28:49}\ts{13}{12}{4:11}. Equivalently, it is the last coefficient $\phi_{kk}$ of the best linear predictor of
$X_t$ from $k$ lags, the coefficient of $X_{t-k}$ in the regression of $X_t$ on $X_{t-1},\dots,X_{t-k}$
\ts{13}{12}{5:13}: the ``incremental correlation'' of adding one more lag. Together the ACF and PACF identify the
order of a model:
\begin{center}\small
\begin{tabular}{@{}lll@{}}
\toprule
 & ACF & PACF\\
\midrule
AR($p$) & decays geometrically or as a damped sinusoid & \textbf{cuts off} after lag $p$\\
MA($q$) & \textbf{cuts off} after lag $q$ & decays\\
ARMA($p,q$) & decays & decays\\
\bottomrule
\end{tabular}
\end{center}
The AR row follows because the regression of an AR($p$) on more than $p$ lags has a zero coefficient on the extra
ones; the MA row is proved in \cref{ch09:sec-ma}.

\begin{casestudy}[10-year Treasury yield: ACF and PACF (2013 Case Study~3)]\label{ch09:cs-3a}
The case study downloads daily constant-maturity yields \texttt{DGS10} from FRED (3\,356 non-missing days, January~2000 to
May~2013; missing values are bond-market holidays), builds weekly (700) and monthly (161) closing series with
\texttt{to.weekly}, \texttt{to.monthly}, and plots ACF and PACF at each frequency. The high first-order
autocorrelation of the levels suggests a unit root at every frequency. The Augmented Dickey--Fuller tests of
\cref{ch09:sec-df} confirm it (p-values 0.057, 0.105 and 0.272 for daily, weekly and monthly levels, none rejecting a unit root
at $5\%$), and reject it for the first differences (p-values $0.01$ at all frequencies, the smallest value R reports).
For the monthly differences $y_0$ (160 values) the table of sample ACF and PACF is
\begin{center}\small
\begin{tabular}{@{}l rrrr@{}}
\toprule
lag & 1 & 2 & 3 & 4\\
\midrule
ACF & $0.035$ & $-0.211$ & $0.079$ & $-0.010$\\
PACF & $0.035$ & $-0.213$ & $0.100$ & $-0.068$\\
\midrule
last coeff.\ of AR($k$) fit by \texttt{lm} & $0.036$ & $-0.217$ & $0.104$ & $-0.072$\\
\bottomrule
\end{tabular}
\end{center}
The last row confirms the definition of the PACF: it is (nearly) the coefficient of the last lag in a regression on $k$ lags.
The small differences arise because \texttt{lm} conditions on the first $k$ cases, whereas \texttt{acf} uses all available
pairs at each lag. Only the lag-2 coefficient is significant ($t=-2.76$ in the AR(2) fit).
\end{casestudy}

\subsubsection{AR(2): roots, correlations and cycles}\label{ch09:sec-ar2}
The AR(2) model $X_t=\phi_0+\phi_1X_{t-1}+\phi_2X_{t-2}+\eta_t$ (2013 PS4 Problem~4, exercise \ref{ch09:ex-ar2})
is the simplest process that can oscillate. Write $\phi(z)=1-\phi_1z-\phi_2z^2=(1-G_1z)(1-G_2z)$, where
$G_i=1/\lambda_i$ are the roots of $G^2-\phi_1G-\phi_2=0$. Then
\[
  G_{1,2}=\frac{\phi_1\pm\sqrt{\phi_1^2+4\phi_2}}2,
\]
and the roots are real and distinct if $\phi_1^2+4\phi_2>0$, real and coincident if $\phi_1^2+4\phi_2=0$, and a
complex-conjugate pair if $\phi_1^2+4\phi_2<0$. Both $|G_i|<1$ if and only if the parameters lie in the
\emph{stationarity triangle}
\[
  \phi_1+\phi_2<1,\qquad \phi_2-\phi_1<1,\qquad |\phi_2|<1 .
\]
The autocorrelations obey $\rho(k)=\phi_1\rho(k-1)+\phi_2\rho(k-2)$ with $\rho(0)=1$ and
$\rho(1)=\phi_1/(1-\phi_2)$ (the first Yule--Walker equation). This is a linear difference equation whose solutions
are combinations of powers of the roots, $\rho(k)=C_1G_1^k+C_2G_2^k$ for $G_1\ne G_2$, and the two initial conditions give
\begin{equation}\label{ch09:eq-ar2rho}
  \rho(k)=\frac{G_1(1-G_2^2)\,G_1^{\,k}-G_2(1-G_1^2)\,G_2^{\,k}}{(G_1-G_2)(1+G_1G_2)} .
\end{equation}
For real roots the correlations decay geometrically in magnitude. For complex roots write
$G_{1,2}=d\,e^{\pm\mathrm i\omega}$ with $d=\sqrt{-\phi_2}$ and $\omega=\arccos\bigl(\phi_1/(2d)\bigr)$; then
\begin{equation}\label{ch09:eq-ar2osc}
  \rho(k)=d^{\,k}\,\frac{\sin(\omega k+F)}{\sin F},\qquad \tan F=\frac{1-\phi_2}{1+\phi_2}\tan\omega,
\end{equation}
a damped cosine with damping factor $d$, frequency $f_0=\omega/2\pi$ (period $1/f_0$) and phase $F$
(taken in the same quadrant as $\omega$). The 2013 problem set states the same result with $|\phi_1|$ and a factor $[\mathrm{sgn}\,\phi_1]^k$. A series that
oscillates around its mean with a typical period of $2\pi/\omega$ periods is the signature of complex roots.

\begin{example}[An AR(2) with a period of eight steps]\label{ch09:ex-ar2osc}
Take $\phi_1=1$, $\phi_2=-0.5$. Then $\phi_1^2+4\phi_2=-1<0$, the roots of $\phi(z)$ are $1\pm\mathrm i$ (modulus
$\sqrt2>1$: stationary), $d=1/\sqrt2$, $\omega=\arccos(1/\sqrt2)=\pi/4$: period exactly $8$, and $\tan F=3$. The
autocorrelations are $\rho=(1,\ 0.667,\ 0.167,\ -0.167,\ -0.25,\ -0.167,\ -0.042,\ 0.042,\dots)$, and the variance is
$\gamma(0)=\sigma^2(1-\phi_2)/[(1+\phi_2)((1-\phi_2)^2-\phi_1^2)]=2.4\,\sigma^2$, equal to $\sigma^2\sum_j\psi_j^2$ for the
Wold weights $\psi=(1,1,0.5,0,-0.25,-0.25,\dots)$ (checked numerically). Note $\psi(1)=1/(1-\phi_1-\phi_2)=2$: a unit shock
eventually raises the level by $2$ (\cref{fig:ar2}).
\end{example}

\begin{figure}[ht]
\centering
\begin{tikzpicture}
\begin{axis}[width=0.95\linewidth,height=5cm,xlabel={lag $k$},ylabel={$\rho(k)$},xmin=-0.5,xmax=20.5,ymin=-0.5,ymax=1.05,grid=major,grid style={black!10}]
  \addplot[dashed,black!50] coordinates {(0.0,1.0) (0.25,0.917) (0.5,0.8409) (0.75,0.7711) (1.0,0.7071) (1.25,0.6484) (1.5,0.5946) (1.75,0.5453) (2.0,0.5) (2.25,0.4585) (2.5,0.4204) (2.75,0.3856) (3.0,0.3536) (3.25,0.3242) (3.5,0.2973) (3.75,0.2726) (4.0,0.25) (4.25,0.2293) (4.5,0.2102) (4.75,0.1928) (5.0,0.1768) (5.25,0.1621) (5.5,0.1487) (5.75,0.1363) (6.0,0.125) (6.25,0.1146) (6.5,0.1051) (6.75,0.0964) (7.0,0.0884) (7.25,0.0811) (7.5,0.0743) (7.75,0.0682) (8.0,0.0625) (8.25,0.0573) (8.5,0.0526) (8.75,0.0482) (9.0,0.0442) (9.25,0.0405) (9.5,0.0372) (9.75,0.0341) (10.0,0.0313) (10.25,0.0287) (10.5,0.0263) (10.75,0.0241) (11.0,0.0221) (11.25,0.0203) (11.5,0.0186) (11.75,0.017) (12.0,0.0156) (12.25,0.0143) (12.5,0.0131) (12.75,0.012) (13.0,0.011) (13.25,0.0101) (13.5,0.0093) (13.75,0.0085) (14.0,0.0078) (14.25,0.0072) (14.5,0.0066) (14.75,0.006) (15.0,0.0055) (15.25,0.0051) (15.5,0.0046) (15.75,0.0043) (16.0,0.0039) (16.25,0.0036) (16.5,0.0033) (16.75,0.003) (17.0,0.0028) (17.25,0.0025) (17.5,0.0023) (17.75,0.0021) (18.0,0.002) (18.25,0.0018) (18.5,0.0016) (18.75,0.0015) (19.0,0.0014) (19.25,0.0013) (19.5,0.0012) (19.75,0.0011) (20.0,0.001)};
  \addplot[dashed,black!50] coordinates {(0.0,-1.0) (0.25,-0.917) (0.5,-0.8409) (0.75,-0.7711) (1.0,-0.7071) (1.25,-0.6484) (1.5,-0.5946) (1.75,-0.5453) (2.0,-0.5) (2.25,-0.4585) (2.5,-0.4204) (2.75,-0.3856) (3.0,-0.3536) (3.25,-0.3242) (3.5,-0.2973) (3.75,-0.2726) (4.0,-0.25) (4.25,-0.2293) (4.5,-0.2102) (4.75,-0.1928) (5.0,-0.1768) (5.25,-0.1621) (5.5,-0.1487) (5.75,-0.1363) (6.0,-0.125) (6.25,-0.1146) (6.5,-0.1051) (6.75,-0.0964) (7.0,-0.0884) (7.25,-0.0811) (7.5,-0.0743) (7.75,-0.0682) (8.0,-0.0625) (8.25,-0.0573) (8.5,-0.0526) (8.75,-0.0482) (9.0,-0.0442) (9.25,-0.0405) (9.5,-0.0372) (9.75,-0.0341) (10.0,-0.0313) (10.25,-0.0287) (10.5,-0.0263) (10.75,-0.0241) (11.0,-0.0221) (11.25,-0.0203) (11.5,-0.0186) (11.75,-0.017) (12.0,-0.0156) (12.25,-0.0143) (12.5,-0.0131) (12.75,-0.012) (13.0,-0.011) (13.25,-0.0101) (13.5,-0.0093) (13.75,-0.0085) (14.0,-0.0078) (14.25,-0.0072) (14.5,-0.0066) (14.75,-0.006) (15.0,-0.0055) (15.25,-0.0051) (15.5,-0.0046) (15.75,-0.0043) (16.0,-0.0039) (16.25,-0.0036) (16.5,-0.0033) (16.75,-0.003) (17.0,-0.0028) (17.25,-0.0025) (17.5,-0.0023) (17.75,-0.0021) (18.0,-0.002) (18.25,-0.0018) (18.5,-0.0016) (18.75,-0.0015) (19.0,-0.0014) (19.25,-0.0013) (19.5,-0.0012) (19.75,-0.0011) (20.0,-0.001)};
  \addplot+[ycomb,thick,mark=*,mark size=1.5pt,color=defcol] coordinates {(0,1.0) (1,0.6667) (2,0.1667) (3,-0.1667) (4,-0.25) (5,-0.1667) (6,-0.0417) (7,0.0417) (8,0.0625) (9,0.0417) (10,0.0104) (11,-0.0104) (12,-0.0156) (13,-0.0104) (14,-0.0026) (15,0.0026) (16,0.0039) (17,0.0026) (18,0.0007) (19,-0.0007) (20,-0.001)};
\end{axis}
\end{tikzpicture}
\caption{Autocorrelation of the AR(2) with $\phi=(1,-0.5)$: a damped cosine of period $8$ under the
envelope $\pm(1/\sqrt2)^k$ (dashed).}\label{fig:ar2}
\end{figure}

\begin{casestudy}[A cycle in the monthly yield changes (2013 Case Study~3)]\label{ch09:cs-3b}
Fitting AR(2) by least squares to the monthly differences of the 10-year yield gives
$\hat\phi_1=0.038$, $\hat\phi_2=-0.217$ (standard errors $0.079$, $0.079$), $R^2=0.048$. The roots of the characteristic
polynomial are the complex pair $0.088\pm2.145\,\mathrm i$, of squared modulus $4.61>1$: the fitted model is
stationary. With complex roots the series is cyclical, with period
$2\pi/\arccos\bigl(|\hat\phi_1|/(2\sqrt{-\hat\phi_2})\bigr)=4.107$, that is, ``just over $4$ months''; I recomputed this value
(and the roots) numerically from the printed coefficients. The function \texttt{ar()} selects the order with the AIC
(\cref{ch09:sec-modelsel}): \texttt{order = 7}. In the AR(7) fit the significant lags are $2$ ($t=-3.10$) and $7$
($t=-2.02$), $R^2=0.123$; the smallest \emph{squared} modulus of the characteristic roots is $1.4027$ (the case study calls it
the ``smallest root modulus'': the modulus is $\sqrt{1.4027}=1.184$), still outside the unit circle. Regression diagnostics
of the AR(7) fit (\texttt{influence.measures}) flag two high-leverage months, October and November 2008, ``the heart of the
financial crisis'' (\cref{ch05:sec-diag}).
\end{casestudy}

\subsection{Moving-average models MA($q$)}\label{ch09:sec-ma}
Set $p=0$: $X_t-\mu=\theta(L)\eta_t=\eta_t+\theta_1\eta_{t-1}+\dots+\theta_q\eta_{t-q}$
\ts{24}{12}{1:14:24}\ts{13}{8}{1:06:12}. Because $\eta$ is white noise, only overlapping terms contribute:
\begin{equation}\label{ch09:eq-ma}
  \begin{gathered}
  \E X_t=\mu,\qquad \gamma(0)=\sigma^2\bigl(1+\theta_1^2+\dots+\theta_q^2\bigr),\\
  \gamma(j)=\begin{cases}\sigma^2\bigl(\theta_j+\theta_{j+1}\theta_1+\dots+\theta_q\theta_{q-j}\bigr),&1\le j\le q,\\ 0,& j>q,\end{cases}
  \end{gathered}
\end{equation}
with $\theta_0=1$ \ts{24}{12}{1:16:26}. The ACF of an MA($q$) \emph{cuts off} at lag $q$ (a complete finite-memory
process: $X_t$ and $X_{t+j}$ share no noise term when $j>q$), which is the identification rule of the table above.

\paragraph{Invertibility.} The process is \emph{invertible} if all roots of $\theta(z)=0$ lie outside the closed unit
disc; then $\theta^{-1}(L)$ exists by \eqref{ch09:eq-geom} and $\eta_t=\theta^{-1}(L)(X_t-\mu)$ is an infinite
autoregression of the observations \ts{24}{12}{1:15:24}. Invertibility matters for two reasons. Identification: for
MA(1), $\rho(1)=\theta/(1+\theta^2)$, which has maximum $1/2$ at $\theta=1$ and is unchanged under
$\theta\mapsto1/\theta$; for instance $\theta=0.25$ and $\theta=4$ both give $\rho(1)=4/17=0.235$, and if $\sigma^2$ is replaced by $\theta^2\sigma^2$ the variance
$\gamma(0)=\sigma^2(1+\theta^2)$ is also unchanged: two different parameter sets with the same second-order structure, hence
indistinguishable from the ACF. Interpretation: only the invertible
representation has $\eta_t$ equal to the one-step forecast error of $X_t$ given its past, which is the Wold innovation; one
therefore always fits the invertible member of the pair.

\subsection{ARMA($p,q$) and the ARMA(1,1) model}\label{ch09:sec-arma11}
A general ARMA($p,q$) is covariance stationary when $\phi(z)$ has all roots outside the unit disc and is invertible when
$\theta(z)$ does. (If $\phi$ and $\theta$ share a root the factors cancel and the model is redundant, a point to remember
when over-fitting.) The Wold weights are the power-series coefficients of $\theta(z)/\phi(z)$, and the
autocovariances satisfy the Yule--Walker recursion $\gamma(j)=\sum_i\phi_i\gamma(j-i)$ for $j>q$.

For ARMA(1,1), $X_t-\phi X_{t-1}=\phi_0+\eta_t+\theta\eta_{t-1}$ (2013 PS4 Problem~6 and 2024 PS4 Problem~6,
exercise \ref{ch09:ex-arma11}), matching the powers in $\psi(L)=(1+\theta L)/(1-\phi L)$ gives
$\psi_0=1$ and $\psi_j=\phi^{j-1}(\phi+\theta)$ for $j\ge1$. Then $\E X_t=\phi_0/(1-\phi)$ and, summing $\sum_j\psi_j^2$,
\begin{equation}\label{ch09:eq-arma11}\begin{gathered}
  \gamma(0)=\sigma^2\,\frac{1+2\phi\theta+\theta^2}{1-\phi^2},\\
  \rho(1)=\frac{(1+\phi\theta)(\phi+\theta)}{1+2\phi\theta+\theta^2}=\phi+\frac{\theta\sigma^2}{\gamma(0)},\\
  \rho(k)=\phi\,\rho(k-1),\ k\ge2 .
\end{gathered}\end{equation}
Compare with the AR(1), where $\rho(k)=\phi^k$ from the very first lag: the ARMA(1,1) also decays geometrically at
rate $\phi$, but from lag~2 on, and the first correlation $\rho(1)$ is a free parameter that need not equal $\phi$. When
$\phi<0$ the ACF of an AR(1) alternates in sign from lag~1; for an ARMA(1,1) it may be positive at lag~1 and alternate
afterwards. Its distinctive pattern is \emph{small first correlation but long memory}: with $\phi=0.9$,
$\theta=-0.85$ one finds $\rho(1)=0.061$ and $\rho(10)=0.024$, while an AR(1) with $\phi=0.9$ has $\rho(1)=0.9$. An economic
index following such a process would show almost no day-to-day autocorrelation together with a slow, persistent
decay. This structure arises for the \emph{squared} residuals of a GARCH(1,1) model, which follow an ARMA(1,1) with
$\phi=\alpha_1+\beta_1$ and $\theta=-\beta_1$ (\cref{ch:volatility}, \ts{13}{11}{5:15}).
The comparison with MA(1) is the reverse: an MA(1) has a single non-zero correlation and then nothing
(2024 PS4 Problem~6(e)).

% ==================================================================
\section{Non-stationarity: differencing, ARIMA and unit-root tests}\label{ch09:sec-nonstat}

Many economic series behave like random walks. Box and Jenkins advocate removing trending behaviour by
\emph{differencing} \ts{24}{12}{1:17:28}\ts{13}{8}{1:07:14}. With $\Delta=1-L$,
\[
  \Delta X_t=X_t-X_{t-1},\qquad
  \Delta^2X_t=(1-L)^2X_t=X_t-2X_{t-1}+X_{t-2},\qquad
  \Delta^k=(1-L)^k=\sum_{j=0}^k\binom kj(-L)^j .
\]
A linear trend $a+bt$ becomes the constant $b$ under $\Delta$; a quadratic trend is removed by $\Delta^2$
\ts{24}{12}{1:18:30}\ts{13}{8}{1:08:20}. In practice one uses $k\le2$. Modelling $\Delta X_t$ or $\Delta^2X_t$ as stationary
means modelling the dynamics of the slope, or of the curvature, of the series, with a constant level and variability
\ts{24}{12}{1:19:32}.

\subsection{Two kinds of trend}\label{ch09:sec-trends}
\paragraph{Trend reversion.} Let $X_t=a+bt+\xi_t$, with $\xi_t=\phi\xi_{t-1}+\eta_t$, $|\phi|<1$
\ts{13}{8}{1:09:21}. Then $\E X_t=a+bt$ (not constant, hence not stationary), $\Var X_t=\sigma^2/(1-\phi^2)$, and
\[
  \Delta X_t=b+\xi_t-\xi_{t-1}=b+(1-L)\xi_t=b+(1-L)(1-\phi L)^{-1}\eta_t,
\]
which is stationary. A shock to $\xi$ dies out and $X$ returns to the deterministic line.

\paragraph{Stochastic trend.} The pure \emph{integrated process} of order one, $I(1)$, is $X_t=X_{t-1}+\eta_t$, that is,
$\Delta X_t=\eta_t$ \ts{13}{8}{1:10:22}. Given $X_0$, $X_t=X_0+TS_t$ with $TS_t=\sum_{j=1}^t\eta_j$ (the
\emph{stochastic trend}), a random walk with white-noise steps, and
\begin{equation}\label{ch09:eq-rwmoments}
  \begin{gathered}
  \Var X_t=t\sigma^2,\qquad \Cov(X_t,X_{t-j})=(t-j)\sigma^2,\\
  \Corr(X_t,X_{t-j})=\frac{t-j}{\sqrt{t(t-j)}}=\sqrt{1-\frac jt}\ \longrightarrow 1\quad(t\to\infty),
  \end{gathered}
\end{equation}
so the process is not stationary: variance and covariances depend on $t$, and consecutive values become
almost perfectly correlated. The trend cannot be forecast (the best forecast at any horizon is the last value, with error
variance growing like $h\sigma^2$).

The two models differ where it matters for finance: in the trend-reversion model the long-run cumulative response
$\psi(1)$ of $\Delta X_t$ to a shock is $\psi(1)=(1-1)/(1-\phi)=0$, the level effect is temporary; for the stochastic trend $\Delta X_t=\eta_t$ has
$\psi(1)=1$, the shock is permanent. Prices are usually modelled as the second kind (\cref{ch04:rem-logprice}),
returns as stationary. Differencing a trend-reversion series (``over-differencing'') creates an MA factor $(1-L)$ with a unit root, which is
not invertible; differencing a stochastic trend is the correct treatment.

\begin{definition}[ARIMA$(p,d,q)$]\label{ch09:def-arima}
$\{X_t\}$ is \emph{integrated ARMA}, ARIMA$(p,d,q)$, if $\Delta^dX_t$ is stationary (and lower-order differencing is not) and
follows an ARMA$(p,q)$ model (2024 slide~25, not discussed in class; \ts{13}{8}{1:11:26}).
\end{definition}
The practical questions (\ts{13}{8}{1:12:26}) are: determine the order $d$ of differencing needed to remove deterministic or
stochastic trends (\cref{ch09:sec-df}); estimate the parameters (\cref{ch09:sec-est}); choose among competing
$(p,d,q)$ (\cref{ch09:sec-modelsel}).

\subsection{Testing for a unit root: the Dickey--Fuller test}\label{ch09:sec-df}
Suppose $X_t=\phi X_{t-1}+\eta_t$ with $\eta_t$ white noise. The \emph{Dickey--Fuller} (DF) test has
\[
  H_0:\ \phi=1\ \text{(unit root, non-stationary)}\qquad\text{against}\qquad H_1:\ |\phi|<1\ \text{(stationary)} .
\]
Fit the AR(1) by least squares and form $t_{\phi=1}=(\hat\phi-1)/\mathrm{se}(\hat\phi)$. Under $H_1$,
$\sqrt T(\hat\phi-\phi)\xrightarrow{d}N(0,1-\phi^2)$ (the standard AR(1) result). Under $H_0$ the estimator is
\emph{superconsistent}, converging at the rate $1/T$, and $T(\hat\phi-1)$, like the $t$ statistic, has a non-standard
distribution, the Dickey--Fuller distribution (2024 slide~28, not discussed in either year). It can be derived in a few lines. For
$X_t=X_{t-1}+\eta_t$ and $X_0=0$,
\[
  \hat\phi-1=\frac{\sum_tX_{t-1}\eta_t}{\sum_tX_{t-1}^2},\qquad
  \sum_{t=1}^TX_{t-1}\eta_t=\tfrac12\Bigl(X_T^2-\sum_{t=1}^T\eta_t^2\Bigr),
\]
because $X_t^2=X_{t-1}^2+2X_{t-1}\eta_t+\eta_t^2$ telescopes. Dividing the numerator by $T\sigma^2$ and the
denominator by $T^2\sigma^2$, and using $X_{\lfloor rT\rfloor}/(\sigma\sqrt T)\to W(r)$ (Donsker's theorem; the
random walk converges to Brownian motion, \cref{ch:sp2}),
\begin{equation}\label{ch09:eq-df}
  T(\hat\phi-1)\ \xrightarrow{d}\ \frac{\tfrac12\bigl(W(1)^2-1\bigr)}{\int_0^1W(r)^2\,\mathrm dr}
  =\frac{\int_0^1W\,\mathrm dW}{\int_0^1W^2\,\mathrm dr},
\end{equation}
the last equality being Itô's formula $\int_0^1W\,\mathrm dW=\tfrac12(W(1)^2-1)$ (\cref{ch:stochcalc}). The distribution is skewed to the
left, so the usual normal or $t$ critical values reject the unit root too often. Simulating $20\,000$ random walks of length $500$,
the $5\%$ critical value of $t_{\phi=1}$ is $-1.95$ (no intercept) and $-2.89$ (with an intercept; the asymptotic value is
about $-2.86$), compared with $-1.645$ for the standard normal; the $1\%$ values are $-2.55$ and $-3.46$.

The test is often \emph{augmented} (ADF, Said--Dickey): lagged differences $\Delta X_{t-1},\dots,\Delta X_{t-k}$ are added to the
regression to absorb serial correlation, and $k$ is chosen by an information criterion. The null hypothesis is
\emph{non-stationarity}, so failing to reject does not prove a unit root: the power is low against $\phi$ near $1$. Other tests:
Phillips--Perron (PP) and Elliott--Rothenberg--Stock (ERS) for a unit-root null, and KPSS, whose null is stationarity
\ts{13}{8}{1:12:26}. The 10-year Treasury yield is the standard illustration (\cref{ch09:cs-3a}): the ADF test does not reject a
unit root for the daily, weekly or monthly \emph{levels} and rejects it for the \emph{differences}. The same test was used in the crack-spread example
(\cref{ch09:sec-crack}).

\subsection{Estimation of ARMA models}\label{ch09:sec-est}
For Gaussian innovations $\eta_t\sim N(0,\sigma^2)$ i.i.d., the parameters are estimated by maximum likelihood
(\cref{ch05:sec-mle}) (2024 slide~26; \ts{13}{8}{1:12:26}).
\begin{itemize}
  \item \emph{AR($p$).} Conditioning on the first $p$ observations, the log-likelihood is
        $\ell=-\frac n2\log(2\pi\sigma^2)-\frac1{2\sigma^2}\sum_{t=p+1}^{T}\bigl(x_t-c-\sum_j\phi_jx_{t-j}\bigr)^2$ with $n=T-p$:
        conditional maximum likelihood \emph{is} least squares of $x_t$ on its $p$ lags (\cref{ch09:sec-wold-recipe}). This is
        the \emph{limited-information} (LIML) method, and the resulting $\hat\sigma^2=\mathrm{RSS}/n$ is
        biased downwards by the same degrees-of-freedom argument as in regression \ts{13}{8}{6:10}\ts{13}{8}{7:11}. The
        \emph{full-information} (FIML) method adds the log-density of the first $p$ values under their stationary distribution,
        which matters when $T$ is small or the roots are close to the unit circle.
  \item \emph{MA and ARMA.} The likelihood is not a regression, because the past innovations are not observed. LIML sets the
        first $q$ innovations to zero and computes the others recursively from the data; the exact likelihood is obtained by
        writing the model in state-space form and applying the Kalman filter (\cref{ch09:sec-ss}), the
        \emph{prediction-error decomposition}. The maximisation is numerical, and stationarity and invertibility constrain the
        parameters; the standard trick is to reparametrise so that the optimiser works on an unconstrained
        scale (logarithms, logits), as in the GARCH fits of \cref{ch:volatility}.
  \item \emph{Method of moments}: Yule--Walker (\cref{ch09:sec-yw}) for AR.
\end{itemize}

\subsection{Model selection: AIC, BIC, HQ}\label{ch09:sec-modelsel}
Fit all ARMA$(p,q)$ with $0\le p\le p_{\max}$, $0\le q\le q_{\max}$, let $\tilde\sigma^2(p,q)$ be the Gaussian MLE of
$\sigma^2$, and choose $(p,q)$ to minimise (2024 slide~27; \ts{13}{8}{1:13:33})
\begin{align}
  \mathrm{AIC}(p,q)&=\log\tilde\sigma^2(p,q)+2\,\frac{p+q}n, &&\text{(Akaike)}\notag\\
  \mathrm{BIC}(p,q)&=\log\tilde\sigma^2(p,q)+\log n\,\frac{p+q}n, &&\text{(Bayes / Schwarz)}\label{ch09:eq-ic}\\
  \mathrm{HQ}(p,q)&=\log\tilde\sigma^2(p,q)+2\log(\log n)\,\frac{p+q}n. &&\text{(Hannan--Quinn)}\notag
\end{align}
Because the maximised Gaussian log-likelihood is $\hat\ell=-\frac n2\bigl(\log(2\pi\tilde\sigma^2)+1\bigr)$, each criterion is
$-2\hat\ell/n$ plus a penalty proportional to the number of parameters (up to an additive constant); the lecture ends on the question
of what penalty is appropriate \ts{13}{8}{1:14:33}. Adding one parameter lowers $\log\tilde\sigma^2$ and raises the penalty. The larger
model is preferred iff $2\Delta\hat\ell>c_n$, the likelihood-ratio statistic (asymptotically $\chi^2_1$ under the smaller model)
exceeds $c_n$, with $c_n=2$ (AIC), $\log n$ (BIC) or $2\log\log n$ (HQ):
\begin{center}\small
\begin{tabular}{@{}l rrr@{}}
\toprule
 & AIC & BIC & HQ\\
\midrule
threshold $c_n$ for $n=100$ & $2$ & $4.61$ & $3.05$\\
\quad equivalent $p$-value & $0.157$ & $0.032$ & $0.081$\\
threshold $c_n$ for $n=500$ & $2$ & $6.21$ & $3.65$\\
\quad equivalent $p$-value & $0.157$ & $0.013$ & $0.056$\\
\bottomrule
\end{tabular}
\end{center}
(For $\chi^2_1$ the AIC threshold $2$ corresponds to $|t|>\sqrt2=1.41$.) So AIC is a liberal criterion, which tends to overfit but
is the right choice if the true model is infinite-dimensional and one wants good forecasts; BIC assumes the true model has a finite number
of parameters that it will find with probability tending to one; HQ makes the same finite-model assumption but with a penalty
growing as slowly as consistency allows \ts{13}{8}{1:15:39}. The same ideas were used for regression in
\cref{ch:regression} and are used in Case Study~3 (\texttt{ar()} minimises the AIC over $p\le22$ and picks $p=7$, with $p=2$ at a
relative AIC of $1.61$) and for VAR models below.

% ==================================================================
\section{Multivariate time series and VAR models}\label{ch09:sec-var}

We now observe $m$ series together, $\bm X_t=(X_{1,t},\dots,X_{m,t})\T\in\R^m$: several exchange rates, a basket of stocks, or
macroeconomic variables \ts{13}{11}{30:44}. $\{\bm X_t\}$ is covariance stationary if every component is
\ts{13}{11}{31:47}.

\subsection{Second-order structure}\label{ch09:sec-multi}
With $\bm\mu=\E\bm X_t$ define the \emph{covariance matrix} $\Gamma_0=\E[(\bm X_t-\bm\mu)(\bm X_t-\bm\mu)\T]$ (symmetric, with the
variances of the components on the diagonal), the \emph{lag-$k$ cross-covariance matrix}
$\Gamma_k=\E[(\bm X_t-\bm\mu)(\bm X_{t-k}-\bm\mu)\T]$, whose $(j,j')$ entry is $\Cov(X_{j,t},X_{j',t-k})$, and the
\emph{correlation matrices} $R_k=D^{-1/2}\Gamma_kD^{-1/2}$ with $D=\mathrm{diag}(\Gamma_0)$
\ts{13}{11}{31:47}\ts{13}{11}{33:03}\ts{13}{11}{35:02}. The matrix $\Gamma_0$ is symmetric, but $\Gamma_k$ for $k\ne0$ is
\emph{not}, and $\Gamma_{-k}=\Gamma_k\T$: the lags of series $j'$ may be correlated with series $j$ without the converse
\ts{13}{11}{36:04}. If $[\Gamma_k]_{jj'}\ne0$ for some $k>0$ one says that $X_{j'}$ \emph{leads} $X_j$ (its past carries information
about $X_j$); if each leads the other there is \emph{feedback} \ts{13}{11}{37:08}\ts{13}{11}{38:13}. This is the idea of
\emph{Granger causality} (Granger and Engle received the Nobel prize for it): lags of one variable help predict another. It is
predictability, not causation in a philosophical sense.

The \emph{multivariate Wold theorem} holds verbatim with matrices:
$\bm X_t=\bm V_t+\sum_{k\ge0}\Psi_k\bm\eta_{t-k}$ with $\Psi_0=I_m$, $\sum_k\Psi_k\Psi_k\T$ convergent, $\bm V_t$ linearly
deterministic, and $\bm\eta_t$ multivariate white noise: $\E\bm\eta_t=\bm 0$, $\Var\bm\eta_t=\Sigma$ (positive semi-definite),
$\Cov(\bm\eta_t,\bm\eta_{t-k})=0$ for $k\ne0$, uncorrelated with $\bm V$ \ts{13}{11}{39:18}\ts{13}{11}{41:24}. The
innovation is $\bm\eta_t=\bm X_t-\E[\bm X_t\mid\mathcal F_{t-1}]$ in the Gaussian case, where $\mathcal F_{t-1}$ is the information available
before time $t$: what is new at $t$ \ts{13}{11}{42:25}\ts{13}{11}{43:32}.

\subsection{VAR($p$) models and their stationarity}\label{ch09:sec-var-def}
\begin{definition}[VAR($p$)]\label{ch09:def-var}
$\{\bm X_t\}$ follows a \emph{vector autoregression} of order $p$ if
\begin{equation}\label{ch09:eq-var}
  \bm X_t=\bm C+\Phi_1\bm X_{t-1}+\Phi_2\bm X_{t-2}+\dots+\Phi_p\bm X_{t-p}+\bm\eta_t,
\end{equation}
with $\bm C\in\R^m$, $\Phi_1,\dots,\Phi_p$ real $m\times m$ matrices and $\{\bm\eta_t\}$ multivariate white noise
$N(\bm 0,\Sigma)$ \ts{13}{11}{45:39}.
\end{definition}
Component $j$ is an AR($p$) of its own lags \emph{plus} lag regressors on all the other series:
$X_{j,t}=c_j+\sum_{k=1}^p[\Phi_k]_{jj}X_{j,t-k}+\sum_{j'\ne j}\sum_{k=1}^p[\Phi_k]_{jj'}X_{j',t-k}+\eta_{j,t}$: exchange rates
that co-move, macro variables that respond to each other \ts{13}{11}{46:43}\ts{13}{11}{47:43}.

\paragraph{VAR(1) form.} Any VAR($p$) is a VAR(1) in higher dimension \ts{13}{11}{48:51}. Stack
\[
  \bm Z_t=(\bm X_t\T,\bm X_{t-1}\T,\dots,\bm X_{t-p+1}\T)\T\in\R^{mp}.
\]
Then
\begin{equation}\label{ch09:eq-companion}
  \bm Z_t=\bm D+A\bm Z_{t-1}+\bm F_t,\qquad
  A=\begin{pmatrix}\Phi_1&\Phi_2&\cdots&\Phi_{p-1}&\Phi_p\\ I_m&0&\cdots&0&0\\ 0&I_m&&0&0\\ \vdots&&\ddots&&\vdots\\ 0&0&\cdots&I_m&0\end{pmatrix},\
  \bm D=\begin{pmatrix}\bm C\\ \bm 0\\ \vdots\\ \bm 0\end{pmatrix},\
  \bm F_t=\begin{pmatrix}\bm\eta_t\\ \bm 0\\ \vdots\\ \bm 0\end{pmatrix},
\end{equation}
with $A$ the ($mp\times mp$) \emph{companion matrix} \ts{13}{11}{50:56}\ts{13}{11}{52:00}.

\begin{proposition}[Stationarity of a VAR]\label{ch09:prop-companion}
The following are equivalent:
\begin{enumerate}[(i)]
  \item all eigenvalues of $A$ have modulus less than~1;
  \item all roots of
      $\det\bigl(I_m-\Phi_1z-\dots-\Phi_pz^p\bigr)=0$ lie outside the closed unit disc $|z|\le1$.
\end{enumerate}
Under (i)--(ii) the VAR has the unique
stationary solution $\bm X_t-\bm\mu=\sum_{k\ge0}\Psi_k\bm\eta_{t-k}$, where $\Psi_k$ is the upper-left $m\times m$ block of $A^k$, and
\begin{equation}\label{ch09:eq-varmean}
\bm\mu=\E\bm X_t=(I_m-\Phi_1-\dots-\Phi_p)^{-1}\bm C .
\end{equation}
\end{proposition}
\begin{proof}
Let $\lambda\ne0$ be an eigenvalue of $A$ with eigenvector $\bm v=(\bm v_1,\dots,\bm v_p)$. The block rows of $A\bm v=\lambda\bm v$
say $\sum_j\Phi_j\bm v_j=\lambda\bm v_1$ and $\bm v_{j-1}=\lambda\bm v_j$ for $j=2,\dots,p$, so $\bm v_j=\lambda^{p-j}\bm v_p$. Substituting,
$\bigl(\lambda^pI-\sum_j\Phi_j\lambda^{p-j}\bigr)\bm v_p=0$ with $\bm v_p\ne0$, i.e.\ $\det\bigl(I-\sum_j\Phi_jz^j\bigr)=0$ at $z=1/\lambda$;
the converse is the same computation backwards. Thus the eigenvalues of $A$ are the reciprocals of the roots, and (i)$\Leftrightarrow$(ii).
If (i) holds, $\bm Z_t-\bm\mu_Z=\sum_{k\ge0}A^k\bm F_{t-k}$ converges because the spectral radius of $A$ is $<1$ (so $\|A^k\|$ decays
geometrically), and reading off the first block gives the expansion. Taking expectations in \eqref{ch09:eq-var} gives $\bm\mu=\bm C+\sum\Phi_j\bm\mu$.
\end{proof}
The characteristic polynomial $\det(I_m-\Phi_1z-\dots-\Phi_pz^p)$ has degree $mp$, hence $mp$ roots
\ts{13}{11}{53:06}\ts{13}{11}{54:08}; the mean-adjusted series satisfies $\bm X_t-\bm\mu=\sum_j\Phi_j(\bm X_{t-j}-\bm\mu)+\bm\eta_t$
\ts{13}{11}{55:13}\ts{13}{11}{56:18}. The scalar case $m=1$ is \cref{ch09:thm-ar-stat}, and its proof-by-eigenvalue is
\cref{ch09:rem-companion}.

\begin{example}[Moments of a VAR(1)]\label{ch09:ex-var1}
For $\bm X_t=\bm C+A\bm X_{t-1}+\bm\eta_t$ with stationary $A$, let $\tilde{\bm X}_t=\bm X_t-\bm\mu$. Then
$\Gamma_0=A\Gamma_0A\T+\Sigma$, and with the vectorisation of \cref{ch09:sec-varest} this is solved by
$\mathrm{vec}\,\Gamma_0=(I-A\otimes A)^{-1}\mathrm{vec}\,\Sigma$; further, $\Gamma_h=A^h\Gamma_0$ for $h\ge0$. Take
$A=\bigl(\begin{smallmatrix}0.5&0.2\\-0.3&0.6\end{smallmatrix}\bigr)$, $\bm C=(1,0.5)\T$,
$\Sigma=\bigl(\begin{smallmatrix}1&0.3\\0.3&0.5\end{smallmatrix}\bigr)$. The eigenvalues of $A$ are $0.55\pm0.24\,\mathrm i$
(modulus $0.6<1$), $\bm\mu=(1.923,-0.192)\T$,
$\Gamma_0=\bigl(\begin{smallmatrix}1.443&0.244\\0.244&0.847\end{smallmatrix}\bigr)$ and
$\Gamma_1=A\Gamma_0=\bigl(\begin{smallmatrix}0.770&0.291\\-0.287&0.435\end{smallmatrix}\bigr)$, which is not symmetric (checked against the
iteration $\Gamma\leftarrow A\Gamma A\T+\Sigma$). This is the pattern that 2013 PS6 Problem~1 asks for with different numbers.
\end{example}

\subsection{VAR as a multivariate regression: estimation}\label{ch09:sec-varest}
Given $n$ observations $\bm x_1,\dots,\bm x_n$ and $p$ pre-sample values, define for each component $j$ the regression
$\bm y^{(j)}=Z\bm\beta^{(j)}+\bm\varepsilon^{(j)}$, where $\bm y^{(j)}=(x_{j,1},\dots,x_{j,n})\T$ and the \emph{same}
$n\times(mp+1)$ design matrix $Z$ has rows $(1,\bm x_{t-1}\T,\dots,\bm x_{t-p}\T)$ \ts{13}{11}{57:18}\ts{13}{11}{58:32}. The $\bm\beta^{(j)}$ contain the
$j$th rows of $\bm C,\Phi_1,\dots,\Phi_p$. Stacking columns,
\begin{equation}\label{ch09:eq-mvreg}
  Y=ZB+E,\qquad Y=[\bm y^{(1)}\cdots\bm y^{(m)}],\ B=[\bm\beta^{(1)}\cdots\bm\beta^{(m)}],
\end{equation}
a multivariate regression, called \emph{seemingly unrelated regressions} (SUR) in econometrics because the $m$ equations are
linked only through the correlation $\Sigma$ of their errors \ts{13}{11}{59:41}\ts{13}{11}{1:01:51}\ts{13}{11}{1:03:53}.

\paragraph{Component-wise least squares.} Fit each equation separately,
$\hat{\bm\beta}^{(j)}=(Z\T Z)^{-1}Z\T\bm y^{(j)}$, that is $\hat B=(Z\T Z)^{-1}Z\T Y$, take the residual matrix
$\hat E=Y-Z\hat B=(I_n-Z(Z\T Z)^{-1}Z\T)Y$, and estimate the innovation covariance by
$\hat\Sigma=\hat E\T\hat E/(n-mp-1)$ (unbiased entry by entry; the lecture notes divide by $n-pm$) \ts{13}{11}{1:04:58}.
Treating the equations separately looks wasteful, since the errors are correlated. It is not:

\begin{theorem}[Optimality of component-wise OLS]\label{ch09:thm-varols}
For a VAR($p$) with unrestricted coefficient matrices,
\begin{enumerate}[(i)]
  \item component-wise OLS equals the generalised least-squares (GLS) estimator for the
      whole system with any innovation covariance $\Sigma$, hence is BLUE (Gauss--Markov, \cref{ch05:thm-gm}, and \cref{ch05:sec-gls});
  \item with Gaussian
      innovations it is the maximum-likelihood estimator of $B$, and the MLE of $\Sigma$ is $\tilde\Sigma=\hat E\T\hat E/n$
      \ts{13}{11}{1:06:58}\ts{13}{11}{1:08:00}.
\end{enumerate}
\end{theorem}
\begin{proof}
Let $\otimes$ be the Kronecker product ($A\otimes B$ is the block matrix with blocks $a_{ij}B$; $(A\otimes B)\T=A\T\otimes B\T$ and
$(A\otimes B)(C\otimes D)=AC\otimes BD$) and $\mathrm{vec}$ the operator that stacks the columns of a matrix into one vector
\ts{13}{11}{1:09:04}\ts{13}{11}{1:10:04}. Put $\bm y=\mathrm{vec}\,Y$, $\mathcal X=I_m\otimes Z$, $\bm\beta^*=\mathrm{vec}\,B$,
$\bm\varepsilon=\mathrm{vec}\,E$. Then $\bm y=\mathcal X\bm\beta^*+\bm\varepsilon$ with $\Cov\bm\varepsilon=\Sigma\otimes I_n$, and
$(\Sigma\otimes I_n)^{-1}=\Sigma^{-1}\otimes I_n$. The GLS estimator is
$[\mathcal X\T(\Sigma^{-1}\otimes I_n)\mathcal X]^{-1}\mathcal X\T(\Sigma^{-1}\otimes I_n)\bm y$ with
\[
  \mathcal X\T(\Sigma^{-1}\otimes I_n)\mathcal X=(I\otimes Z\T)(\Sigma^{-1}\otimes I_n)(I\otimes Z)=\Sigma^{-1}\otimes Z\T Z,\qquad
  \mathcal X\T(\Sigma^{-1}\otimes I_n)=\Sigma^{-1}\otimes Z\T ,
\]
so it equals $[\Sigma\otimes(Z\T Z)^{-1}](\Sigma^{-1}\otimes Z\T)\bm y=[I_m\otimes(Z\T Z)^{-1}Z\T]\bm y=\mathrm{vec}\bigl((Z\T Z)^{-1}Z\T Y\bigr)$:
$\Sigma$ cancels and one recovers component-wise OLS \ts{13}{11}{1:11:06}\ts{13}{11}{1:12:08}. This proves (i).
For (ii), the Gaussian log-likelihood is
\[
  \ell(B,\Sigma)=-\tfrac{nm}2\log2\pi-\tfrac n2\log|\Sigma|-\tfrac12\mathrm{tr}\bigl[\Sigma^{-1}(Y-ZB)\T(Y-ZB)\bigr]
\]
\ts{13}{11}{1:13:12}\ts{13}{11}{1:14:12}. Writing $Y-ZB=\hat E+Z(\hat B-B)$ and using $Z\T\hat E=0$, the trace equals
$\mathrm{tr}[\Sigma^{-1}\hat E\T\hat E]+\mathrm{tr}[\Sigma^{-1}(\hat B-B)\T Z\T Z(\hat B-B)]$, and the second term is $\ge0$; so for \emph{every}
$\Sigma$ the maximum over $B$ is at $\hat B$ (a \emph{concentrated} likelihood, because $\hat B$ does not depend on $\Sigma$
\ts{13}{11}{1:15:16}\ts{13}{11}{1:16:19}). The remaining function is $-\frac n2[\log|\Sigma|+\mathrm{tr}(\Sigma^{-1}S)]$ with
$S=\hat E\T\hat E/n$. Let $u_1,\dots,u_m>0$ be the eigenvalues of $\Sigma^{-1/2}S\Sigma^{-1/2}$; then
$\log|\Sigma|+\mathrm{tr}(\Sigma^{-1}S)=\log|S|+\sum_i(u_i-\log u_i)\ge\log|S|+m$, since $u-\log u\ge1$, with equality iff all $u_i=1$,
i.e.\ $\Sigma=S$. (This is the Anderson--Olkin result quoted in the lecture \ts{13}{11}{1:16:19}.)
\end{proof}

\paragraph{Order selection and asymptotics.} As in the univariate case, fit $\mathrm{VAR}(p)$ for $0\le p\le p_{\max}$ and choose $p$ by
\begin{equation}\label{ch09:eq-icvar}
  \begin{gathered}
  \mathrm{AIC}(p)=\log|\tilde\Sigma(p)|+2\,\frac{pm^2}n,\qquad
  \mathrm{BIC}(p)=\log|\tilde\Sigma(p)|+\log n\,\frac{pm^2}n,\\
  \mathrm{HQ}(p)=\log|\tilde\Sigma(p)|+2\log(\log n)\,\frac{pm^2}n
  \end{gathered}
\end{equation}
\ts{13}{11}{1:17:24}. (The slides print $-\log|\tilde\Sigma(p)|$, a sign misprint: $\log|\tilde\Sigma|$ is the log
generalised variance, which decreases as $p$ increases, while the penalty increases; compare \cref{ch09:eq-ic}.) For a stationary VAR whose
innovations have bounded fourth moments and are independent over time, $Q=\mathrm{plim}\,Z\T Z/n$ exists and is non-singular and
$\sqrt n(\hat{\bm\beta}^*-\bm\beta^*)\to N(0,\Sigma\otimes Q^{-1})$; least-squares and Gaussian maximum-likelihood estimates are asymptotically
identical (2013 notes, last slide).

\subsection{Impulse responses and a macroeconomic VAR}\label{ch09:sec-irf}
The matrices $\Psi_k$ of the Wold expansion are the \emph{impulse response functions}: $[\Psi_k]_{jj'}$ is the effect on variable $j$, $k$ periods
later, of a unit innovation in variable $j'$ \ts{13}{11}{1:21:42}\ts{13}{11}{1:22:42}\ts{13}{11}{1:23:43}. Because $\Sigma$ is not
diagonal the innovations are correlated, and R's \texttt{irf} plots the \emph{orthogonalised} responses (shocks made uncorrelated through a Cholesky
factor of $\Sigma$, which depends on the ordering of the variables; a point the course does not discuss).

\begin{casestudy}[VAR of US macroeconomic series (2013 Case Study~5 and 2013 L11--L12)]\label{ch09:cs-5}
In the 1980s Litterman and Sims used vector autoregressions of nine macroeconomic aggregates (Treasury-bill rate, M1, GNP deflator, real GNP,
business fixed investment, unemployment, trade-weighted dollar, S\&P~500, commodity price index) to forecast the economy and to assess
the effect of policy variables such as the Federal Funds rate. The case study downloads monthly series from FRED with
\texttt{quantmod} (UNRATE, FEDFUNDS, TB3MS, CPIAUCSL, M1SL, GDPDEF, GDP, GPDI, TWEXBMTH, and the S\&P~500 from Yahoo) and fits VAR models with the
\texttt{vars} package to three of them: the unemployment rate, the Federal Funds rate and the CPI, over 1960--2000 (492 months)
\ts{13}{11}{1:18:29}\ts{13}{11}{1:19:33}\ts{13}{11}{1:20:38}.
\begin{itemize}
  \item \emph{Levels.} \texttt{VARselect} (up to 12 lags) returns $p=12,5,2,12$ for AIC, HQ, SC (=BIC), FPE. The VAR(2) chosen by SC has
        characteristic-root moduli $1.002,\,0.986,\,0.952,\,0.468,\,0.331,\,0.084$: the largest exceeds 1, so the fitted model is
        \emph{non-stationary} \ts{13}{12}{0:00}\ts{13}{12}{1:08}. (Here ``roots'' are eigenvalues of the companion matrix, stationary when $<1$;
        \cref{ch09:rem-companion}.)
  \item \emph{Differences.} Removing the random-walk component by first differences, the ACF and PACF show strong self-correlation and some cross-correlation between
        unemployment and Fed funds \ts{13}{12}{2:11}\ts{13}{12}{3:11}\ts{13}{12}{4:11}. \texttt{VARselect} gives $p=12,3,3,12$; the VAR(3) has
        all root moduli $<0.84$ (stationary). The dependent variable is the monthly change; the lecture's reading of the Fed funds equation
        \ts{13}{12}{6:15}\ts{13}{12}{7:21}:
        \[
          \Delta\mathrm{FF}_t=-0.71\,\Delta\mathrm{UN}_{t-1}+0.37\,\Delta\mathrm{FF}_{t-1}-0.18\,\Delta\mathrm{FF}_{t-2}+\dots\quad(t=-4.9,\ 8.0,\ -3.7).
        \]
        If unemployment rises, the Fed lowers rates the next period; a rate change tends to be followed by another in the same direction
        (the Fed does not want to shock the economy), and since $0.37\,x_{t-1}-0.18\,x_{t-2}=0.19\,x_{t-1}+0.18\,(x_{t-1}-x_{t-2})$ the
        negative second lag is an \emph{acceleration} term \ts{13}{12}{8:21}\ts{13}{12}{9:24}\ts{13}{12}{10:30}. The CPI equation is dominated by its own lags
        ($0.42$, $t=9.5$, and $0.29$ at lag~3, $t=6.7$), the unemployment equation by lags 2 and 3 of itself; $R^2$ are $0.12$ (unemployment),
        $0.23$ (Fed funds), $0.49$ (CPI); the residual correlation between unemployment and Fed funds is $-0.20$.
  \item \emph{Impulse responses} (bootstrap $95\%$ bands, 100 runs) \ts{13}{11}{1:22:42}. Levels VAR: after a rise in unemployment the Fed funds rate
        and the CPI decline (less employment, less pressure on prices); after a rise in the Fed funds rate unemployment and the CPI increase; after a
        rise in the CPI the Fed funds rate increases (consistent with policy against inflation). Differenced VAR: unemployment up $\Rightarrow$ Fed funds
        down over the next quarters; Fed funds up $\Rightarrow$ a modest increase in inflation, which persists for several quarters (the large third-lag
        partial autocorrelation of the CPI); CPI up $\Rightarrow$ modestly higher unemployment and higher Fed funds.
\end{itemize}
Three variables are far too few for a structural model of the economy: the lecture recommends ten to twelve
\ts{13}{12}{5:13}.
\end{casestudy}

% ==================================================================
\section{Cointegration}\label{ch09:sec-coint}

Differencing removes non-stationarity but discards information about the \emph{levels}. Cointegration is the framework that keeps it
\ts{13}{12}{10:30}\ts{13}{12}{11:36}.

\subsection{Definition and examples}
A process $\{\bm X_t\}$ is \emph{integrated of order $d$}, $I(d)$, if $(1-L)^d\bm X_t$ is stationary. If it has a VAR representation
$\Phi(L)\bm X_t=\bm\eta_t$ then $\Phi(L)=(1-L)^d\Phi^*(L)$ with $\Phi^*$ describing the stationary VAR($p-d$) of $\Delta^d\bm X_t$
\ts{13}{12}{12:43}\ts{13}{12}{13:43}. The point is that all the components may be $I(1)$ without the vector being ``jointly'' $I(1)$: linear
combinations of the levels may be stationary \ts{13}{12}{14:47}.

\begin{definition}[Cointegration]\label{ch09:def-coint}
An $m$-vector $\{\bm X_t\}$ of $I(1)$ series is \emph{cointegrated} if there is $\bm\beta\in\R^m$, $\bm\beta\ne\bm0$, with $\bm\beta\T\bm X_t$ an $I(0)$ (stationary)
process \ts{13}{12}{15:49}. The vector $\bm\beta$ is defined up to scale; normalising $\beta_1=1$,
$\bm\beta\T\bm X_t=u_t$ says $x_{1,t}=-(\beta_2x_{2,t}+\dots+\beta_mx_{m,t})+u_t$: a \emph{long-run equilibrium relationship}
$x_1=-(\beta_2x_2+\dots)$ with a stationary \emph{disequilibrium error} (cointegration residual) $u_t$ \ts{13}{12}{16:49}.
\end{definition}

\begin{finance}[Where cointegration appears]\label{ch09:fin-coint}
Examples listed in the lecture \ts{13}{12}{16:49}\ts{13}{12}{17:50}: the term structure of interest rates (the level of rates is non-stationary, the
\emph{spreads} between maturities are stationary: expectations hypothesis); purchasing-power parity (exchange rate and price levels of two countries: otherwise there would be arbitrage);
money demand (money, income, prices, interest rates); covered interest-rate parity (spot and forward rates); the law of one price (identical or equivalent
assets that must be valued alike: spot and futures, the same asset on two venues). In trading, a cointegrated pair is the basis of
\emph{convergence} (pairs) trades: two prices that wander individually but whose spread mean-reverts.
\end{finance}

\subsection{The error-correction (VECM) representation}\label{ch09:sec-vecm}
Let $\bm X_t$ follow a VAR($p$) \eqref{ch09:eq-var} and be $I(1)$. Subtract $\bm X_{t-1}$ from both sides, then add and subtract
$(\Phi_1-I)\bm X_{t-2}$, then $(\Phi_2+\Phi_1-I)\bm X_{t-3}$, and so on \ts{13}{12}{27:08}\ts{13}{12}{28:12}\ts{13}{12}{29:13}. The result is
\begin{equation}\label{ch09:eq-vecm}
  \begin{gathered}
  \Delta\bm X_t=\bm C+\Pi\,\bm X_{t-1}+\Gamma_1\Delta\bm X_{t-1}+\dots+\Gamma_{p-1}\Delta\bm X_{t-p+1}+\bm\eta_t,\\
  \Pi=\Phi_1+\dots+\Phi_p-I_m,\qquad \Gamma_k=-\sum_{j=k+1}^p\Phi_j,
  \end{gathered}
\end{equation}
the \emph{vector error-correction model} (VECM) \ts{13}{12}{30:14}\ts{13}{12}{31:14}. (Proof: for $p=2$, $\bm X_t=\bm C+\Phi_1\bm X_{t-1}+\Phi_2\bm X_{t-2}+\bm\eta_t$ is written
$\bm X_{t-2}=\bm X_{t-1}-\Delta\bm X_{t-1}$ in the last term, giving $\Delta\bm X_t=\bm C+(\Phi_1+\Phi_2-I)\bm X_{t-1}-\Phi_2\Delta\bm X_{t-1}+\bm\eta_t$; the general case
follows by induction.) Everything in \eqref{ch09:eq-vecm} is $I(0)$ except possibly $\Pi\bm X_{t-1}$, which must therefore be stationary too: it contains the
cointegrating combinations. Since $\bm X$ has unit roots, $\Pi$ is singular. Three cases \ts{13}{12}{32:19}\ts{13}{12}{33:19}:
\begin{itemize}
  \item $\mathrm{rank}\,\Pi=0$, $\Pi=0$: no cointegration; model the differences by a stationary VAR($p-1$);
  \item $\mathrm{rank}\,\Pi=m$: $\bm X$ is itself stationary (no unit roots);
  \item $0<r=\mathrm{rank}\,\Pi<m$: $\Pi=\alpha\beta\T$ with $\alpha,\beta$ of size $m\times r$ and full rank $r$. The $r$ columns of $\beta$ are linearly independent
        cointegrating vectors; the matrix $\alpha$ holds the \emph{adjustment} (error-correction) coefficients: the process is pulled back towards the equilibrium
        $\beta\T\bm X=0$ at speed $\alpha$.
\end{itemize}
The factorisation is not unique: $\alpha\beta\T=(\alpha G)(\beta G^{-\mathsf T})\T$ for any invertible $r\times r$ matrix $G$. The stationary
directions form an $r$-dimensional space; a choice of coordinates in it is convention \ts{13}{12}{34:22}.

\begin{example}[A cointegrated pair]\label{ch09:ex-coint}
Let $x_t=x_{t-1}+\varepsilon_{1t}$ (random walk), $u_t=0.7u_{t-1}+\varepsilon_{2t}$ and $y_t=2x_t+u_t$. Both $x$ and $y$ are $I(1)$, but $y_t-2x_t=u_t$ is
stationary: $\bm\beta=(-2,1)\T$. In error-correction form
$\Delta x_t=\varepsilon_{1t}$ and $\Delta y_t=2\Delta x_t+u_t-u_{t-1}=-0.3\,(y_{t-1}-2x_{t-1})+2\varepsilon_{1t}+\varepsilon_{2t}$. Hence
$\Delta\bm X_t=\alpha\,\bm\beta\T\bm X_{t-1}+\text{noise}$ with $\alpha=(0,-0.3)\T$, and $\Pi=\alpha\bm\beta\T=\bigl(\begin{smallmatrix}0&0\\0.6&-0.3\end{smallmatrix}\bigr)$
has rank one. In the VAR(1) in levels the matrix $I+\Pi$ has eigenvalues $1$ (the common stochastic trend) and $0.7$ (the stationary spread).
Only $y$ reacts to a disequilibrium: $x$ is ``weakly exogenous'' (the row of $\alpha$ for $x$ is zero). \Cref{ch09:fig-coint} shows a simulated path.
\end{example}

\begin{figure}[ht]
\centering
\begin{tikzpicture}
\begin{axis}[width=0.95\linewidth,height=5.4cm,xlabel={$t$},xmin=0,xmax=119,grid=major,grid style={black!10},legend style={font=\footnotesize,at={(0.5,1.03)},anchor=south,legend columns=3}]
  \addplot[thick,color=defcol] coordinates {(0,0.0) (1,0.3) (2,0.03) (3,-0.86) (4,-1.32) (5,-2.31) (6,-2.25) (7,-0.91) (8,-1.4) (9,-2.02) (10,-1.53) (11,-1.18) (12,-1.07) (13,-2.0) (14,-2.03) (15,-1.34) (16,-2.68) (17,-3.14) (18,-5.04) (19,-6.33) (20,-8.17) (21,-8.41) (22,-9.67) (23,-9.4) (24,-9.24) (25,-9.43) (26,-11.95) (27,-12.49) (28,-12.54) (29,-12.42) (30,-13.95) (31,-14.43) (32,-15.41) (33,-16.22) (34,-15.16) (35,-15.96) (36,-16.0) (37,-15.11) (38,-15.7) (39,-15.81) (40,-15.7) (41,-15.63) (42,-16.86) (43,-16.78) (44,-15.42) (45,-16.97) (46,-16.11) (47,-15.99) (48,-16.63) (49,-14.63) (50,-13.87) (51,-15.07) (52,-15.0) (53,-14.42) (54,-14.61) (55,-13.92) (56,-13.99) (57,-13.32) (58,-11.89) (59,-12.56) (60,-12.36) (61,-12.82) (62,-12.69) (63,-13.88) (64,-14.46) (65,-14.66) (66,-13.76) (67,-12.61) (68,-13.94) (69,-14.73) (70,-14.08) (71,-16.08) (72,-16.54) (73,-16.64) (74,-15.38) (75,-14.69) (76,-15.02) (77,-15.39) (78,-15.64) (79,-14.11) (80,-14.54) (81,-14.84) (82,-14.49) (83,-14.61) (84,-14.81) (85,-15.92) (86,-15.94) (87,-16.38) (88,-15.21) (89,-14.56) (90,-14.58) (91,-13.92) (92,-14.26) (93,-13.2) (94,-13.21) (95,-12.63) (96,-13.92) (97,-13.57) (98,-15.26) (99,-17.29) (100,-17.6) (101,-18.5) (102,-18.33) (103,-16.09) (104,-16.92) (105,-17.54) (106,-17.34) (107,-16.85) (108,-17.02) (109,-17.23) (110,-16.53) (111,-16.01) (112,-17.04) (113,-17.12) (114,-17.08) (115,-18.14) (116,-17.88) (117,-18.74) (118,-17.76) (119,-17.57)};
  \addplot[thick,color=thmcol] coordinates {(0,0.0) (1,0.65) (2,-0.27) (3,-2.02) (4,-4.04) (5,-6.28) (6,-5.45) (7,-3.76) (8,-3.66) (9,-5.69) (10,-3.76) (11,-3.35) (12,-2.37) (13,-4.08) (14,-5.04) (15,-2.61) (16,-4.45) (17,-5.68) (18,-9.82) (19,-12.57) (20,-16.87) (21,-16.52) (22,-19.47) (23,-18.92) (24,-19.05) (25,-19.63) (26,-25.2) (27,-25.13) (28,-25.27) (29,-24.41) (30,-27.59) (31,-29.06) (32,-31.15) (33,-33.0) (34,-30.71) (35,-32.43) (36,-32.52) (37,-31.42) (38,-32.71) (39,-31.55) (40,-31.75) (41,-32.15) (42,-34.13) (43,-33.01) (44,-31.33) (45,-34.41) (46,-32.93) (47,-33.53) (48,-33.91) (49,-29.73) (50,-28.02) (51,-30.79) (52,-30.17) (53,-29.29) (54,-29.62) (55,-28.79) (56,-29.37) (57,-26.82) (58,-24.2) (59,-25.24) (60,-24.82) (61,-25.98) (62,-25.93) (63,-27.76) (64,-29.1) (65,-29.53) (66,-27.65) (67,-24.62) (68,-27.04) (69,-28.65) (70,-27.94) (71,-32.82) (72,-32.98) (73,-32.62) (74,-30.39) (75,-28.79) (76,-29.16) (77,-29.66) (78,-29.94) (79,-27.57) (80,-27.71) (81,-29.48) (82,-28.32) (83,-28.46) (84,-28.56) (85,-29.98) (86,-29.67) (87,-31.91) (88,-30.84) (89,-28.92) (90,-29.64) (91,-28.17) (92,-28.24) (93,-27.21) (94,-28.24) (95,-26.37) (96,-28.59) (97,-27.82) (98,-30.97) (99,-35.42) (100,-36.69) (101,-38.14) (102,-38.05) (103,-34.13) (104,-34.91) (105,-35.87) (106,-34.98) (107,-34.5) (108,-35.0) (109,-35.73) (110,-34.47) (111,-32.89) (112,-35.16) (113,-34.78) (114,-34.34) (115,-35.19) (116,-35.83) (117,-36.99) (118,-35.24) (119,-34.95)};
  \addplot[thick,color=intcol] coordinates {(0,0.0) (1,0.05) (2,-0.32) (3,-0.29) (4,-1.4) (5,-1.66) (6,-0.95) (7,-1.94) (8,-0.85) (9,-1.64) (10,-0.7) (11,-0.99) (12,-0.23) (13,-0.08) (14,-0.98) (15,0.06) (16,0.91) (17,0.6) (18,0.25) (19,0.08) (20,-0.53) (21,0.29) (22,-0.12) (23,-0.12) (24,-0.56) (25,-0.77) (26,-1.3) (27,-0.16) (28,-0.2) (29,0.44) (30,0.31) (31,-0.2) (32,-0.33) (33,-0.57) (34,-0.39) (35,-0.5) (36,-0.53) (37,-1.2) (38,-1.32) (39,0.07) (40,-0.36) (41,-0.88) (42,-0.41) (43,0.55) (44,-0.48) (45,-0.46) (46,-0.7) (47,-1.55) (48,-0.64) (49,-0.46) (50,-0.28) (51,-0.65) (52,-0.18) (53,-0.45) (54,-0.4) (55,-0.95) (56,-1.39) (57,-0.17) (58,-0.43) (59,-0.12) (60,-0.11) (61,-0.34) (62,-0.54) (63,-0.0) (64,-0.18) (65,-0.22) (66,-0.14) (67,0.61) (68,0.83) (69,0.81) (70,0.23) (71,-0.67) (72,0.1) (73,0.65) (74,0.37) (75,0.59) (76,0.88) (77,1.11) (78,1.33) (79,0.66) (80,1.37) (81,0.21) (82,0.67) (83,0.76) (84,1.06) (85,1.87) (86,2.2) (87,0.85) (88,-0.42) (89,0.2) (90,-0.47) (91,-0.34) (92,0.27) (93,-0.8) (94,-1.82) (95,-1.12) (96,-0.76) (97,-0.68) (98,-0.45) (99,-0.83) (100,-1.49) (101,-1.14) (102,-1.38) (103,-1.95) (104,-1.06) (105,-0.78) (106,-0.3) (107,-0.81) (108,-0.96) (109,-1.27) (110,-1.42) (111,-0.88) (112,-1.08) (113,-0.55) (114,-0.18) (115,1.09) (116,-0.07) (117,0.48) (118,0.28) (119,0.19)};
  \legend{$x_t$ (random walk),$y_t=2x_t+u_t$,$u_t=y_t-2x_t$ (stationary)}
\end{axis}
\end{tikzpicture}
\caption{A simulated cointegrated pair (seed~7). Each of $x$ and $y$ wanders like a random walk, but $y-2x$ stays around zero.}\label{ch09:fig-coint}
\end{figure}

\subsection{Estimation and tests for the cointegration rank}\label{ch09:sec-cointest}
Sims, Stock and Watson (1990) and Park and Phillips (1989) prove that least squares applied to the original VAR in \emph{levels}
gives consistent estimates with the same asymptotic distributions as maximum likelihood, and the reduced-rank constraints on $\Pi$ hold
asymptotically: no new tool is needed for estimation \ts{13}{12}{34:22}\ts{13}{12}{35:33}\ts{13}{12}{36:40}. For inference on the \emph{rank}
$r$, Johansen (1991) developed maximum likelihood by reduced-rank regression for Gaussian innovations, following Banerjee and Hendry's use of the
concentrated likelihood \ts{13}{12}{37:43}. It gives likelihood-ratio tests of the number of cointegrating vectors: Johansen's \emph{trace} statistic
$-T\sum_{i>r}\log(1-\hat\lambda_i)$ (a sum over the smallest eigenvalues, testing rank $\le r$) and his \emph{maximum-eigenvalue} statistic
$-T\log(1-\hat\lambda_{r+1})$ (rank $r$ against $r+1$), where the $\hat\lambda_i$ are squared sample canonical correlations between
$\Delta\bm X_t$ and $\bm X_{t-1}$ after adjusting for the lagged differences (the slides describe them as the eigenvalues of $\hat\Pi$; I give the standard
definition). A cruder route, not in the course, is the Engle--Granger two-step method: regress one series on the others and apply an ADF test
to the residual (with different critical values).

\subsection{Energy futures and crack spreads}\label{ch09:sec-crack}
The lecture's example of cointegration uses futures on the CME from the Energy Information Administration: WTI crude oil (CL, quoted in dollars per barrel), RBOB
gasoline and heating oil (both in cents per gallon), from 2006 to 2013, first-month contracts. A refinery ``cracks'' crude oil into gasoline, heating oil, jet fuel;
multiplying the product prices by $42$ gallons per barrel puts inputs and outputs in the same units, and the products stay above crude, so the
\emph{crack spreads} (product minus crude) represent the value added of refining \ts{13}{12}{18:54}\ts{13}{12}{19:55}\ts{13}{12}{20:56}\ts{13}{12}{22:01}. The class
listed what drives the spread: cost of refining, supply and demand, seasonality (heating oil dearer in winter), and the refiner's profit
\ts{13}{12}{23:02}\ts{13}{12}{24:05}; the spreads look stationary \ts{13}{12}{25:05}. The case note (not in the download) tests the levels with the Augmented Dickey--Fuller test:
p-values $0.164$ for the first crude contract and $0.327$ for heating oil, so a unit root cannot be rejected \ts{13}{12}{38:47}\ts{13}{12}{39:49}. In the Johansen
procedure the test of rank $0$ is almost significant at the $10\%$ level and the test of rank $\le1$ is not, so the rank is at most one, with cointegrating vector $(1,\,1.3,\,-1.7)$ on crude, RBOB and heating oil, so that
``crude plus gasoline minus heating oil'' is the stationary combination (ignoring seasonality) \ts{13}{12}{41:55}\ts{13}{12}{43:00}. The first-month contract was chosen because
refiners hedge on the futures market, so the analysis addresses hedging refinery risk; longer maturities would be used depending on what one hedges \ts{13}{12}{44:07}.
The video jumps here between explanations, so the numbers are reported as spoken; commodity models are the subject of \cref{ch:commodities}.

% ==================================================================
\section{Linear state-space models and the Kalman filter}\label{ch09:sec-ss}

Many models of economics and finance, including ARMA models, regressions with time-varying coefficients and models with hidden factors, can be written in
one common form that has an exact and cheap recursive solution, the Kalman filter \ts{13}{12}{45:11}\ts{13}{12}{46:12}.

\subsection{The model}\label{ch09:sec-ss-model}
Let $\bm y_t\in\R^k$ be the observation at time $t$, $\bm s_t\in\R^m$ the (hidden) \emph{state}, $\bm\varepsilon_t\in\R^k$ the observation error and
$\bm\eta_t\in\R^n$ the state innovation. The \emph{linear Gaussian state-space model} is
\begin{equation}\label{ch09:eq-ss}
  \bm s_{t+1}=T_t\bm s_t+R_t\bm\eta_t,\quad \bm\eta_t\sim N(\bm0,Q_t);\qquad
  \bm y_t=Z_t\bm s_t+\bm\varepsilon_t,\quad \bm\varepsilon_t\sim N(\bm0,H_t),
\end{equation}
the \emph{state} (transition) equation and the \emph{observation} (measurement) equation, with i.i.d.\ noises independent of one another;
$T_t$ ($m\times m$) is the transition matrix, $R_t$ ($m\times n$, often columns of the identity) selects which components are shocked, $Z_t$ ($k\times m$) the observation
matrix, $Q_t,H_t$ positive definite \ts{13}{12}{47:18}\ts{13}{12}{48:22}. Jointly,
$\bigl(\bm s_{t+1};\bm y_t\bigr)=\bigl(\begin{smallmatrix}T_t\\ Z_t\end{smallmatrix}\bigr)\bm s_t+\bm u_t$ with
$\bm u_t=(R_t\bm\eta_t;\bm\varepsilon_t)\sim N(\bm0,\Omega_t)$, $\Omega_t=\mathrm{blockdiag}(R_tQ_tR_t\T,\,H_t)$ \ts{13}{12}{49:25}\ts{13}{12}{50:28}. Often the
model is time-invariant. The state is not unique: replacing $\bm s_t$ by $M\bm s_t$ ($M$ invertible) and $T_t$ by $MT_tM^{-1}$, $Z_t$ by $Z_tM^{-1}$ gives an equivalent model
\ts{13}{12}{1:02:59}.

\subsection{Examples}\label{ch09:sec-ss-ex}
\paragraph{Time-varying CAPM.} Extend the CAPM regression (\cref{ch05:capm}) so that alpha and beta drift as Gaussian random walks:
\[
  r_t=\alpha_t+\beta_tr_{m,t}+\varepsilon_t,\quad \varepsilon_t\sim N(0,\sigma_\varepsilon^2);\qquad
  \alpha_{t+1}=\alpha_t+\xi_t,\quad\beta_{t+1}=\beta_t+\zeta_t,
\]
with $\xi_t\sim N(0,\sigma_\xi^2)$, $\zeta_t\sim N(0,\sigma_\zeta^2)$ and $\varepsilon,\xi,\zeta$ mutually independent, $r_t$ the excess return of the asset and $r_{m,t}$ that of the market
\ts{13}{12}{48:22}\ts{13}{12}{51:33}. In state-space form $\bm s_t=(\alpha_t,\beta_t)\T$, $T_t=R_t=I_2$, $Q=\mathrm{diag}(\sigma_\xi^2,\sigma_\zeta^2)$,
$Z_t=(1,\ r_{m,t})$ (a known $1\times2$ row), $H=\sigma_\varepsilon^2$ \ts{13}{12}{52:35}\ts{13}{12}{53:36}.

\paragraph{Regression with time-varying coefficients.} $y_t=\bm x_t\T\bm\beta_t+\varepsilon_t$ with $\beta_{j,t+1}=\beta_{j,t}+\eta_{j,t}$, $\eta_{j,t}\sim N(0,\sigma_j^2)$, is
$\bm s_t=\bm\beta_t$, $T_t=R_t=I_p$, $Z_t=\bm x_t\T$, $Q=\mathrm{diag}(\sigma_1^2,\dots,\sigma_p^2)$, $H=\sigma^2$ \ts{13}{12}{54:37}. If all $\sigma_j^2=0$ the coefficients are constant
and one obtains ordinary linear regression; applying the filter of the next subsection observation by observation is then \emph{recursive least squares}, updating the estimate as
data arrive \ts{13}{12}{55:37}.

\paragraph{AR($p$).} With $\bm s_t=(y_t,y_{t-1},\dots,y_{t-p+1})\T$,
$\bm s_{t+1}=F\bm s_t+\bm e_1\eta_t$, where $F$ is the companion matrix of \cref{ch09:rem-companion}, $\bm e_1=(1,0,\dots,0)\T$, and $y_t=\bm e_1\T\bm s_t$ (no measurement
error, $\varepsilon_t=0$) \ts{13}{12}{56:43}\ts{13}{12}{57:46}\ts{13}{12}{58:46}.

\paragraph{ARMA($p,q$) (Harvey's representation).} Older estimation methods for moving-average models were hard to implement; the state-space form made exact maximum likelihood routine
\ts{13}{12}{59:48}\ts{13}{12}{1:00:55}. Put $m=\max(p,q+1)$, $\phi_j=0$ for $j>p$, $\theta_j=0$ for $j>q$, $\theta_0=1$; then $\{y_t\}$ is ARMA($m,m-1$). Define the state by $s_{1,t}=y_t$ and
\[
  s_{k,t}=\sum_{i=k}^{m}\phi_i\,y_{t+k-1-i}+\sum_{j=k-1}^{m-1}\theta_j\,\eta_{t+k-1-j},\qquad k=2,\dots,m .
\]
Rewriting the model equation in terms of $s_{1,t}$ and $s_{2,t}$ and repeating gives
\begin{equation}\label{ch09:eq-harvey}
  \bm s_{t+1}=T\bm s_t+\bm R\,\eta_{t+1},\quad y_t=\bm e_1\T\bm s_t,\qquad
  T=\begin{pmatrix}\phi_1&1&0&\cdots&0\\ \phi_2&0&1&&\vdots\\ \vdots&&&\ddots&0\\ \phi_{m-1}&0&\cdots&0&1\\ \phi_m&0&\cdots&0&0\end{pmatrix},\quad
  \bm R=\begin{pmatrix}1\\ \theta_1\\ \vdots\\ \theta_{m-1}\end{pmatrix}
\end{equation}
\ts{13}{12}{1:01:55}\ts{13}{12}{1:06:01}. \emph{Check for ARMA(1,1)} ($m=2$): the first row of $\bm s_{t+1}=T\bm s_t+\bm R\eta_{t+1}$ is
$y_{t+1}=\phi y_t+s_{2,t}+\eta_{t+1}$ and the second is $s_{2,t+1}=\theta\eta_{t+1}$, so $s_{2,t}=\theta\eta_t$ and $y_{t+1}=\phi y_t+\eta_{t+1}+\theta\eta_t$, as required. I also
verified numerically that for ARMA(2,1) with $\phi=(0.5,-0.3)$, $\theta=0.4$ the state-space variances ($\gamma(0)=1.893$, $\gamma(1)=1.036$ in units of $\sigma^2$, from the discrete Lyapunov equation
$P=TPT\T+\bm R\bm R\T$) agree with the Wold-weight sums. The MA($q$) is the case $p=0$. (The lecture notes give a different MA($q$) form, with the state $(\eta_{t-1},\dots,\eta_{t-q})$ and
$y_t=\sum_j\theta_j s_j+\eta_t$; there the measurement error and the state innovation are the \emph{same} variable, which violates the independence assumption of \eqref{ch09:eq-ss}, so I use the ARMA form above.)

\subsection{The Kalman filter}\label{ch09:sec-kalman}
Let $\mathcal F_t=\{\bm y_1,\dots,\bm y_t\}$. The \emph{Kalman filter} is the recursive computation of the conditional laws of the state and of the next observation given the information so far
\ts{13}{12}{1:07:03}\ts{13}{12}{1:08:10}. Because everything is Gaussian these are determined by means and covariances:
\begin{align*}
  \bm s_{t|t}&=\E[\bm s_t\mid\mathcal F_t], & \bm s_{t|t-1}&=\E[\bm s_t\mid\mathcal F_{t-1}], & \bm y_{t|t-1}&=\E[\bm y_t\mid\mathcal F_{t-1}],\\
  P_{t|t}&=\Cov(\bm s_t\mid\mathcal F_t), & P_{t|t-1}&=\Cov(\bm s_t\mid\mathcal F_{t-1}), & F_t&=\Cov(\bm y_t\mid\mathcal F_{t-1}),
\end{align*}
where the lecture writes $\Omega_s(t\mid t)$, $\Omega_s(t\mid t-1)$, $\Omega_y(t\mid t-1)$ \ts{13}{12}{1:09:12}. The \emph{innovation} (one-step forecast error) is
$\bm v_t=\bm y_t-\bm y_{t|t-1}=\bm y_t-Z_t\bm s_{t|t-1}$ \ts{13}{12}{1:10:12}. The filter alternates two steps.
\begin{enumerate}[1.]
  \item \textbf{Prediction} \ts{13}{12}{1:11:18}. From $(\bm s_{t|t},P_{t|t})$ predict the next state and observation:
        \begin{equation}\label{ch09:eq-kfpred}
          \begin{gathered}
          \bm s_{t+1|t}=T_t\bm s_{t|t},\qquad P_{t+1|t}=T_tP_{t|t}T_t\T+R_tQ_tR_t\T,\\
          \bm y_{t+1|t}=Z_{t+1}\bm s_{t+1|t},\qquad F_{t+1}=Z_{t+1}P_{t+1|t}Z_{t+1}\T+H_{t+1}.
          \end{gathered}
        \end{equation}
  \item \textbf{Correction (filtering)} \ts{13}{12}{1:12:21}. When $\bm y_t$ is observed, update:
        \begin{equation}\label{ch09:eq-kfupd}
          \bm s_{t|t}=\bm s_{t|t-1}+G_t\bm v_t,\qquad P_{t|t}=P_{t|t-1}-G_tF_tG_t\T,\qquad G_t=P_{t|t-1}Z_t\T F_t^{-1},
        \end{equation}
        $G_t$ being the \emph{filter gain}. The uncertainty on the state decreases by the information in the innovation (\ts{13}{12}{1:13:23}). (The slides write the gain with the index
        $(t-1\mid t)$; it must be $(t\mid t-1)$.)
\end{enumerate}
Two further operations reuse the same quantities. \emph{Forecasting} $h$ steps ahead iterates the prediction step without corrections \ts{13}{12}{1:14:24}. \emph{Smoothing} improves the estimate of a past
state using the whole sample $\mathcal F_T$, by the backward recursion
$\bm s_{t|T}=\bm s_{t|t}+S_t(\bm s_{t+1|T}-\bm s_{t+1|t})$ with $S_t=P_{t|t}T_t\T P_{t+1|t}^{-1}$ (and the analogous recursion for the covariance) \ts{13}{12}{1:14:24}.

\begin{proof}[Derivation of the correction step]
Conditionally on $\mathcal F_{t-1}$ the pair $(\bm s_t,\bm y_t)$ is jointly Gaussian with means $(\bm s_{t|t-1},\,Z_t\bm s_{t|t-1})$ and covariance blocks
$\Var\bm s_t=P_{t|t-1}$, $\Cov(\bm s_t,\bm y_t)=P_{t|t-1}Z_t\T$ (because $\bm\varepsilon_t$ is independent of $\bm s_t$), $\Var\bm y_t=F_t=Z_tP_{t|t-1}Z_t\T+H_t$. Conditioning a Gaussian on part of
its coordinates (\cref{ch03:sec-mvn}) gives the mean $\bm s_{t|t-1}+P_{t|t-1}Z_t\T F_t^{-1}(\bm y_t-Z_t\bm s_{t|t-1})$ and the covariance
$P_{t|t-1}-P_{t|t-1}Z_t\T F_t^{-1}Z_tP_{t|t-1}=P_{t|t-1}-G_tF_tG_t\T$, which are \eqref{ch09:eq-kfupd}. The prediction step is the linear transformation of a Gaussian by the
state equation.
\end{proof}

\paragraph{Likelihood.} By the chain rule $p(\bm y_1,\dots,\bm y_T;\theta)=p(\bm y_1;\theta)\prod_{t\ge2}p(\bm y_t\mid\mathcal F_{t-1};\theta)$, and
$\bm y_t\mid\mathcal F_{t-1}\sim N(\bm y_{t|t-1},F_t)$. So the filter delivers the exact log-likelihood of all parameters $\theta$ in $T_t,Z_t,R_t,Q_t,H_t$ as a by-product
(\emph{prediction-error decomposition}) \ts{13}{12}{1:15:25}:
\begin{equation}\label{ch09:eq-kfll}
  \ell(\theta)=-\frac{kT}2\log2\pi-\frac12\sum_{t=1}^T\log|F_t|-\frac12\sum_{t=1}^T\bm v_t\T F_t^{-1}\bm v_t .
\end{equation}
It is maximised by Newton-type methods or by the EM algorithm of Dempster, Laird and Rubin (\cref{ch07:sec-fa}). As $T\to\infty$ the MLE is consistent and asymptotically normal
with covariance the inverse Fisher information (2013 notes, last slide). The initial state is modelled as $N(\bm s_{1|0},P_{1|0})$ with pre-specified moments: a large variance (diffuse) if one knows nothing, or the
stationary distribution of the state if the model is stationary \ts{13}{12}{1:16:25}. The observations are jointly Gaussian and the recursions provide a full specification of the model: means and
covariances characterise the whole law.

\begin{example}[The local-level model: Kalman filter, EWMA and ARIMA(0,1,1)]\label{ch09:ex-locallevel}
The simplest model with a hidden state is $y_t=\mu_t+\varepsilon_t$, $\mu_{t+1}=\mu_t+\eta_t$: a slowly moving true level observed with noise. Here $T=Z=R=1$,
$Q=\Var\eta$, $H=\Var\varepsilon$. The recursions \eqref{ch09:eq-kfpred}--\eqref{ch09:eq-kfupd} become $P_{t+1|t}=P_{t|t}+Q$, $G_t=P_{t|t-1}/(P_{t|t-1}+H)$ and
$m_t=(1-G_t)\,m_{t-1}+G_t\,y_t$, where $m_t=\mu_{t|t}$. At the steady state $P=P_{t|t-1}$ solves $P=(1-K)P+Q$ with $K=P/(P+H)$, that is $P^2-QP-QH=0$:
\[
  P=\frac{Q+\sqrt{Q^2+4QH}}2,\qquad K=\frac P{P+H},\qquad P_{t|t}=(1-K)P .
\]
For $Q=1$, $H=4$: $P=(1+\sqrt{17})/2=2.562$, $K=0.390$, $P_{t|t}=1.562$, while the raw observation has error variance $H=4$
(a Monte Carlo over $2\,000$ paths gives a filter mean-squared error of $1.558$). The filter is an \emph{exponentially weighted moving average} (EWMA) of the observations with
weight $K$ and decay factor $1-K=0.610$. The same object appears in ARIMA language: $\Delta y_t=\eta_{t-1}+\varepsilon_t-\varepsilon_{t-1}$ has $\gamma(0)=Q+2H$, $\gamma(1)=-H$, so
$\rho(1)=-H/(Q+2H)=-4/9$, an MA(1) whose invertible parameter is $\theta=-0.6096$ (from $\theta/(1+\theta^2)=\rho(1)$), equal to $-(1-K)$. Thus the local-level model \emph{is} ARIMA(0,1,1),
and the steady-state Kalman filter is the exponentially weighted forecast (\cref{ch09:sec-arma11,ch09:sec-nonstat}). The same recursion applied to squared returns is the RiskMetrics
volatility estimator, a GARCH(1,1) with $\alpha_0=0$ and $\alpha_1+\beta_1=1$ (\cref{ch:volatility}) \ts{13}{11}{27:33}. \Cref{ch09:fig-kf} shows a simulated path.
\end{example}

\begin{figure}[ht]
\centering
\begin{tikzpicture}
\begin{axis}[width=0.95\linewidth,height=5.4cm,xlabel={$t$},xmin=0,xmax=59,grid=major,grid style={black!10},legend style={font=\footnotesize,at={(0.5,1.03)},anchor=south,legend columns=3}]
  \addplot[only marks,mark=*,mark size=1.2pt,color=black!45] coordinates {(0,-0.47) (1,2.96) (2,1.74) (3,2.07) (4,2.5) (5,-0.47) (6,3.05) (7,1.59) (8,3.53) (9,-0.35) (10,4.44) (11,3.38) (12,2.49) (13,3.97) (14,4.85) (15,6.38) (16,0.47) (17,5.18) (18,4.19) (19,3.76) (20,1.07) (21,4.08) (22,-0.56) (23,2.42) (24,1.97) (25,-4.64) (26,-2.52) (27,-5.72) (28,-4.11) (29,-1.53) (30,0.38) (31,-7.45) (32,-5.79) (33,-2.85) (34,-0.38) (35,-3.0) (36,-4.79) (37,-1.66) (38,-2.49) (39,-3.68) (40,-4.39) (41,-1.9) (42,-0.96) (43,-3.05) (44,-3.97) (45,-6.93) (46,-7.07) (47,-3.56) (48,-5.82) (49,-9.06) (50,-2.13) (51,-2.91) (52,0.87) (53,-2.82) (54,-2.91) (55,-6.22) (56,-0.06) (57,3.03) (58,-5.51) (59,-4.82)};
  \addplot[thick,color=defcol] coordinates {(0,0.03) (1,1.39) (2,2.62) (3,2.11) (4,1.81) (5,1.28) (6,1.85) (7,1.8) (8,2.54) (9,0.7) (10,2.26) (11,2.17) (12,2.85) (13,2.71) (14,2.33) (15,2.79) (16,3.62) (17,3.42) (18,3.26) (19,3.95) (20,3.08) (21,1.56) (22,1.96) (23,1.29) (24,-0.63) (25,-1.45) (26,-1.91) (27,-3.11) (28,-4.6) (29,-4.56) (30,-3.67) (31,-3.9) (32,-4.64) (33,-4.26) (34,-3.54) (35,-3.84) (36,-3.29) (37,-2.25) (38,-2.46) (39,-3.27) (40,-2.92) (41,-2.68) (42,-1.58) (43,-2.86) (44,-3.52) (45,-4.36) (46,-6.1) (47,-5.97) (48,-5.44) (49,-6.18) (50,-4.8) (51,-3.97) (52,-3.35) (53,-2.94) (54,-1.99) (55,-3.32) (56,-2.71) (57,-2.1) (58,-3.87) (59,-3.53)};
  \addplot[thick,color=thmcol] coordinates {(0,-0.33) (1,1.28) (2,1.48) (3,1.72) (4,2.02) (5,1.05) (6,1.83) (7,1.73) (8,2.44) (9,1.35) (10,2.55) (11,2.88) (12,2.72) (13,3.21) (14,3.85) (15,4.84) (16,3.13) (17,3.93) (18,4.03) (19,3.93) (20,2.81) (21,3.31) (22,1.79) (23,2.04) (24,2.01) (25,-0.59) (26,-1.34) (27,-3.05) (28,-3.47) (29,-2.71) (30,-1.5) (31,-3.83) (32,-4.59) (33,-3.91) (34,-2.53) (35,-2.71) (36,-3.52) (37,-2.8) (38,-2.68) (39,-3.07) (40,-3.59) (41,-2.93) (42,-2.16) (43,-2.51) (44,-3.08) (45,-4.58) (46,-5.55) (47,-4.77) (48,-5.18) (49,-6.7) (50,-4.91) (51,-4.13) (52,-2.18) (53,-2.43) (54,-2.62) (55,-4.02) (56,-2.48) (57,-0.32) (58,-2.35) (59,-3.31)};
  \legend{observations $y_t$,true level $\mu_t$,filtered $\mu_{t|t}$}
\end{axis}
\end{tikzpicture}
\caption{Local-level model with $Q=1$, $H=4$ (seed~11). The filter (red) tracks the hidden level (blue) much more closely than the raw observations (dots).}\label{ch09:fig-kf}
\end{figure}

\begin{finance}[Where state-space models appear]
Time-varying betas and alphas (the CAPM example above) and, in general, any regression whose coefficients are believed to drift; dynamic factor models, where the factors of
\cref{ch:pca} follow their own time-series model; a ``true'' price observed through noisy quotes, as in the USD/RUB forward-point example of \cref{ch:markets}, where five
brokers quote at random and a Kalman filter estimates the underlying price \ts{13}{1}{57:41}; and exact maximum-likelihood estimation of ARMA models
(\cref{ch09:sec-est}). The filter is cheap: each step costs a few matrix products, and it delivers the likelihood as well as the forecast.
\end{finance}

% ==================================================================
\section{Notes on the sources}\label{ch09:sec-notes}

\begin{coursenote}[2013 versus 2024, and errata]
\begin{itemize}
  \item \textbf{Overlap.} 2024 L12 and 2013 L8 use the same slides for univariate theory (\file{lec11-12} vs.\ \file{lecnote8}); the 2013 lecture goes further (trend reversion, ARIMA,
        estimation, model selection, with the information-criteria discussion at the end). Multivariate series, cointegration, state space and the Kalman filter are 2013-only.
  \item \textbf{Slide misprints.}
  \begin{itemize}
    \item Yule--Walker slide: the second lag is written $X_{t-1}$ instead of $X_{t-2}$ (corrected aloud, \ts{24}{12}{1:12:16}), and the sum in the estimator of $\sigma^2$
        starts at $j-1$ instead of $j=1$.
    \item 2013 L11, model selection for VAR: the criteria are printed with $-\log|\tilde\Sigma(p)|$; it should be $+\log|\tilde\Sigma(p)|$ (\cref{ch09:eq-icvar}).
    \item 2013 L12, Kalman filter: the gain $G_t$ is written with covariances at $(t-1\mid t)$; it should be $(t\mid t-1)$.
    \item The notes divide $\tilde\Sigma$ for the VAR by $n-pm$; the unbiased
        divisor is $n-mp-1$ when an intercept is estimated.
  \end{itemize}
  \item \textbf{Case Study 3.} The printed ``smallest root modulus $1.4027$'' is the smallest \emph{squared} modulus (\cref{ch09:cs-3b}). The AR(2) cycle of ``just over four months'' is
        reproduced ($4.107$).
  \item \textbf{Conventions.} $\mathrm{Var}(\hat r_k)\approx1/(T-k)$ (2024) versus $1/T$ (Bartlett): equivalent for $T\gg k$. R's \texttt{vars} reports moduli of companion eigenvalues, whereas
        \texttt{polyroot} returns roots of $\phi(z)$ (\cref{ch09:rem-companion}). ``Log returns'' of yields is a loose name for relative changes of the yield.
  \item \textbf{Not in the download.} The plots \file{TimeSeries4plots.pdf}, \file{TimeSeries4acfplots.pdf} and the crack-spread cointegration note; the numbers of \cref{ch09:sec-crack} are as spoken.
\end{itemize}
\end{coursenote}

\begin{intuition}[How the pieces fit together]
The chapter is one idea seen from four sides. Stationarity plus the Wold theorem say that a stationary series is a linear filter of white noise. ARMA models parametrise that filter with few
parameters; the AR part is a regression on the past (so least squares of \cref{ch:regression} applies) and the roots of its polynomial decide stability. Non-stationary series are
brought to that setting by differencing, and unit-root tests decide when. In several dimensions the polynomial becomes a matrix (VAR), stability is a condition on the eigenvalues of
the companion matrix, and cointegration is the case where a rank-deficient matrix $\Pi$ ties non-stationary series together. State-space models recast all of it as a hidden Markov-type
Gaussian system, in which the Kalman filter is the exact recursive regression on the past. Forward links: volatility (GARCH is ARMA for squared returns, \cref{ch:volatility}), portfolio
theory (\cref{ch:portfolio}), continuous time (the AR(1) is the Ornstein--Uhlenbeck process, \cref{ch:sde}; the Dickey--Fuller limit uses Brownian motion, \cref{ch:sp2}).
\end{intuition}

% ==================================================================
\begin{keyformulas}
\begin{align*}
&\text{Stationarity:} && \E X_t=\mu,\quad \gamma(\tau)=\Cov(X_t,X_{t+\tau}),\quad\rho=\gamma/\gamma(0)\\
&\text{White noise test:} && \mathrm{BP}=T\textstyle\sum_{j\le K}\hat r_j^2\sim\chi^2_K,\quad |\hat r_k|>1.96/\sqrt T\\
&\text{Wold:} && X_t=V_t+\psi(L)\eta_t,\quad \psi_0=1,\ \textstyle\sum\psi_j^2<\infty\\
&\text{ARMA:} && \phi(L)(X_t-\mu)=\theta(L)\eta_t,\quad \psi(L)=\theta(L)/\phi(L)\\
&\text{Stationary if:} && \text{roots of }\phi(z)\text{ outside }|z|\le1\\
& && (\Leftrightarrow\ \text{eigenvalues of companion matrix }|\lambda|<1)\\
&\text{AR(1):} && \gamma(0)=\sigma^2/(1-\phi^2),\quad \rho(j)=\phi^j\\
&\text{Yule--Walker:} && \gamma(j)=\textstyle\sum_{i\le p}\phi_i\gamma(j-i)+\sigma^2\delta_{j0}\\
&\text{MA($q$):} && \gamma(0)=\sigma^2\textstyle\sum_{i\le q}\theta_i^2,\quad \gamma(j)=0\ (j>q)\\
&\text{ARMA(1,1):} && \gamma(0)=\sigma^2\tfrac{1+2\phi\theta+\theta^2}{1-\phi^2},\quad \rho(k)=\phi\rho(k-1)\ (k\ge2)\\
&\text{Differencing:} && \Delta=1-L;\ \text{ARIMA}(p,d,q):\ \Delta^dX\sim\text{ARMA}(p,q)\\
&\text{Dickey--Fuller:} && T(\hat\phi-1)\to\tfrac{\frac12(W(1)^2-1)}{\int_0^1W^2\dd r}\ \text{under }\phi=1\\
&\text{Information criteria:} && \log\tilde\sigma^2+c_n\tfrac{p+q}n,\quad c_n=2,\ \log n,\ 2\log\log n\\
&\text{VAR($p$):} && \bm X_t=\bm C+\textstyle\sum_j\Phi_j\bm X_{t-j}+\bm\eta_t,\quad \bm\mu=(I-\textstyle\sum\Phi_j)^{-1}\bm C\\
&\text{VAR(1) moments:} && \Gamma_0=A\Gamma_0A\T+\Sigma,\quad \Gamma_h=A^h\Gamma_0\\
&\text{Estimation:} && \hat B=(Z\T Z)^{-1}Z\T Y,\quad \tilde\Sigma=\hat E\T\hat E/n\\
&\text{VECM:} && \Delta\bm X_t=\bm C+\Pi\bm X_{t-1}+\textstyle\sum_k\Gamma_k\Delta\bm X_{t-k}+\bm\eta_t\\
& && \Pi=\textstyle\sum_j\Phi_j-I=\alpha\beta\T,\quad \mathrm{rank}\,\Pi=r\\
&\text{State space:} && \bm s_{t+1}=T\bm s_t+R\bm\eta_t,\quad \bm y_t=Z\bm s_t+\bm\varepsilon_t\\
&\text{Kalman:} && \bm s_{t|t}=\bm s_{t|t-1}+G_t\bm v_t,\quad G_t=P_{t|t-1}Z\T F_t^{-1}\\
& && P_{t|t}=P_{t|t-1}-G_tF_tG_t\T,\quad P_{t+1|t}=TP_{t|t}T\T+RQR\T
\end{align*}
\end{keyformulas}

% ==================================================================
\section*{Exercises}
\addcontentsline{toc}{section}{Exercises}

\subsection*{From 2013 Problem Set 4 (and 2024 Problem Set 4, Problems 5--6)}
The 2013 PS4 has six problems on univariate models; 2024 PS4 Problem~5 is 2013 Problem~1(a)--(e), and 2024 Problem~6 is 2013 Problem~6 plus a part (e).

\begin{exercise}[PS4 (2013), Problem 1; PS4 (2024), Problem 5: covariance-stationary AR(2)]\label{ch09:ex-ps4-1}
Let $X_t=\phi_0+\phi_1X_{t-1}+\phi_2X_{t-2}+\eta_t$, $\eta_t\sim\mathrm{WN}(0,\sigma^2)$, be covariance stationary.
\begin{enumerate}[(a)]
  \item Show that $\mu=\E X_t=\phi_0/(1-\phi_1-\phi_2)$.
  \item Show that $\gamma(k)=\phi_1\gamma(k-1)+\phi_2\gamma(k-2)$, $k\ge1$.
  \item Show that $\rho_k=\phi_1\rho_{k-1}+\phi_2\rho_{k-2}$.
  \item The Yule--Walker equations for $k=1,2$ are $\rho_1=\phi_1+\phi_2\rho_1$ and $\rho_2=\phi_1\rho_1+\phi_2$ (using $\rho_0=1$, $\rho_{-k}=\rho_k$): solve for $\rho_1,\rho_2$ in terms of $\phi_1,\phi_2$.
  \item Solve for $\phi_1,\phi_2$ in terms of $\rho_1,\rho_2$.
  \item (2013 only) Derive complete formulas for $\rho_k$, $k>2$, using Problem~4(b) below.
\end{enumerate}
\end{exercise}

\begin{exercise}[PS4 (2013), Problem 2: difference equations of an AR($p$)]\label{ch09:ex-ps4-2}
Let $\phi(L)X_t=\phi_0+\eta_t$ be an AR($p$) with $\phi(L)=1-\phi_1L-\dots-\phi_pL^p$, and consider the homogeneous equation $\phi(L)g(t)=0$.
\begin{enumerate}[(a)]
  \item Show $\mu=\phi_0/(1-\sum_i\phi_i)$, so $\mu=0$ iff $\phi_0=0$
      (``de-meaning'' removes the constant).
  \item For $\gamma(t)=\Cov(X_t,X_0)$ prove $\phi(L)\gamma(t)=0$ for all $t\ge1$.
  \item Same for $\rho(t)=\Corr(X_t,X_0)$.
  \item Given $x_0,\dots,x_t$, let $\hat x_t(h)=\E[X_{t+h}\mid x_t,\dots,x_0]$ for $h>0$
      and $\hat x_t(h)=x_{t+h}$ for $h\le0$; show $\hat x_t(h)=\phi_0+\sum_i\phi_i\hat x_t(h-i)$ for every $h>0$.
  \item For the de-meaned process show that the forecast function satisfies $\phi(L)\hat x_t(h)=0$ in $h$.
\end{enumerate}
\end{exercise}

\begin{exercise}[PS4 (2013), Problem 3: solutions of the difference equations]\label{ch09:ex-ps4-3}
Consider $\phi(L)g(t)=0$, $t=0,1,2,\dots$
\begin{enumerate}[(a)]
  \item Show that $g(t)=Ce^{-\lambda t}$ solves it if $z=e^{\lambda}$ is a root of $\phi(z)=1-\sum_i\phi_iz^i=0$.
  \item Show that $g(t)=CG^t$ solves it if $G^{-1}$ is a root.
  \item If the $p$ roots
      $z_i=G_i^{-1}$ are distinct, show that $g(t)=\sum_iC_iG_i^t$ is a solution.
  \item If all roots are outside the unit circle, show every such solution is bounded; since the autocovariance is a solution, this is the necessary and sufficient
      condition for covariance stationarity.
\end{enumerate}
\end{exercise}

\begin{exercise}[PS4 (2013), Problem 4: the AR(2) process]\label{ch09:ex-ar2}
\leavevmode
\begin{enumerate}[(a)]
  \item Solve for the two roots $z_1,z_2$ of the characteristic equation and give the conditions under which they are distinct and real, coincident and real, or complex conjugates.
  \item Prove $\rho_k=C_1G_1^k+C_2G_2^k$
      with $G_i=z_i^{-1}$; solve for $C_1,C_2$ when the roots are distinct and show that $\rho_k=\dfrac{G_1(1-G_2^2)G_1^k-G_2(1-G_1^2)G_2^k}{(G_1-G_2)(1+G_1G_2)}$.
  \item For distinct real roots show that the autocovariance decreases geometrically in magnitude.
  \item Under what conditions on $\phi_1$ is it always positive, and when does it alternate in sign?
  \item \emph{Optional.} For complex roots write $G_{1,2}=de^{\pm\mathrm i2\pi f_0}$ and show
      $\rho_k=[\mathrm{sgn}(\phi_1)]^k|d|^k\sin(2\pi f_0k+F)/\sin F$ with $|d|=\sqrt{-\phi_2}$, $2\pi f_0=\arccos\bigl(|\phi_1|/(2\sqrt{-\phi_2})\bigr)$ and $\tan F=\frac{1+d^2}{1-d^2}\tan(2\pi f_0)$.
\end{enumerate}
\end{exercise}

\begin{exercise}[PS4 (2013), Problem 5: the MA(1) process]\label{ch09:ex-ps4-5}
$X_t=\eta_t+\theta\eta_{t-1}$.
\begin{enumerate}[(a)]
  \item Derive the ACF.
  \item Compute the first-order autocorrelation for $\theta=0.5$ and $\theta=2.0$.
  \item Solve the formula for $\rho_1$ for $\theta$ in terms of $\rho_1$; is the solution unique?
  \item For each of the two models decide whether the process is invertible and, if so, write it as an infinite-order autoregression.
\end{enumerate}
\end{exercise}

\begin{exercise}[PS4 (2013), Problem 6; PS4 (2024), Problem 6: ARMA(1,1)]\label{ch09:ex-arma11}
For the covariance-stationary $(1-\phi L)X_t=\phi_0+(1+\theta L)\eta_t$:
\begin{enumerate}[(a)]
  \item Prove $\mu=\phi_0/(1-\phi)$.
  \item Prove $\Var X_t=\gamma(0)=\sigma^2\bigl[1+\frac{(\phi+\theta)^2}{1-\phi^2}\bigr]$
      (equivalently $\sigma^2(1+2\phi\theta+\theta^2)/(1-\phi^2)$).
  \item Prove that $\rho_1=\phi+\theta\sigma^2/\gamma(0)$ and $\rho_k=\phi\rho_{k-1}$ for $k>1$.
  \item Compare the ACF with that of the AR(1) with the same $\phi_0,\phi$: which pattern of the ACF is impossible for an AR(1)?
      What behaviour would an economic index following such a process show?
  \item (2024 only) Compare with the ACF of an MA(1) with the same parameters; what is the pattern of an MA(1) and how does it change for MA($q$), $q>1$?
      Try to obtain (b)--(c) directly from the covariance recursion rather than via the Wold weights used in \cref{ch09:sec-arma11}.
\end{enumerate}
\end{exercise}

\subsection*{From 2013 Problem Set 6 (Problems 1--2)}
\begin{exercise}[PS6 (2013), Problem 1: moments of a VAR(1)]\label{ch09:ex-ps6-1}
Let $\bm X_t=(X_{1,t},X_{2,t})\T$ follow $X_{1,t}=0.3+0.8X_{1,t-1}+\eta_{1,t}$, $X_{2,t}=0.2+0.6X_{1,t-1}+0.4X_{2,t-1}+\eta_{2,t}$, with $(\eta_{1,t},\eta_{2,t})\T$ i.i.d.\ $N(\bm0,\Sigma)$, $\Sigma=\bigl(\begin{smallmatrix}3&-1\\-1&3\end{smallmatrix}\bigr)$.
Compute
\begin{enumerate}[(a)]
  \item $\bm\mu=\E\bm X_t$,
  \item $\Gamma(0)=\Cov\bm X_t$,
  \item $\Gamma(1)=\Cov(\bm X_t,\bm X_{t-1})$, and
  \item derive a formula for $\Gamma(h)$, $h\ge1$. (See \cref{ch09:ex-var1} for the method.)
\end{enumerate}
\end{exercise}

\begin{exercise}[PS6 (2013), Problem 2: an AR(2) as a difference of two MA processes]\label{ch09:ex-ps6-2}
With $\{\eta_t\}$ i.i.d.\ $\mathrm{WN}(0,\sigma^2)$ let $w_t=5(1-0.5L)^{-1}\eta_t$, $v_t=4(1-0.4L)^{-1}\eta_t$ and $x_t=w_t-v_t$.
\begin{enumerate}[(a)]
  \item Solve for the coefficients $\psi_i$ in $x_t=\eta_t+\sum_{i\ge1}\psi_i\eta_{t-i}$.
  \item Prove $\{x_t\}$ is an AR(2) process.
  \item Solve for $\phi_1,\phi_2$ in $x_t=\phi_1x_{t-1}+\phi_2x_{t-2}+\eta_t$.
  \item Prove that any stationary AR(2) process can be written as the difference of two (possibly infinite-order) moving-average processes driven by the same $\{\eta_t\}$.
\end{enumerate}
\end{exercise}

\subsection*{Extra exercises}
\begin{exercise}[Box--Pierce by hand (not from the course)]\label{ch09:ex-x-bp}
A daily return series with $T=500$ has $\hat r_1=0.09$, $\hat r_2=-0.05$, $\hat r_3=0.04$, $\hat r_4=0.02$, $\hat r_5=-0.06$.
\begin{enumerate}[(a)]
  \item Which lags are individually significant at $5\%$?
  \item Compute BP for $K=5$ and compare with $\chi^2_5$ (critical value $11.07$).
  \item Explain how the two answers can disagree and which you would report.
  \item If the numbers came from the residuals of an ARMA(1,1) fit, what would the degrees of freedom be?
\end{enumerate}
\end{exercise}

\begin{exercise}[Stationarity of a VAR (not from the course)]\label{ch09:ex-x-companion}
\leavevmode
\begin{enumerate}[(a)]
  \item Write the companion matrix of the univariate AR(2) $X_t=0.5X_{t-1}+0.3X_{t-2}+\eta_t$, find its eigenvalues ($0.852,-0.352$) and the roots of $\phi(z)$; verify $\lambda=1/z$.
  \item For the VAR(1) with $A=\bigl(\begin{smallmatrix}0.9&0.4\\0&0.5\end{smallmatrix}\bigr)$ compute $\Gamma_0$ from $\mathrm{vec}\,\Gamma_0=(I-A\otimes A)^{-1}\mathrm{vec}\,\Sigma$ for $\Sigma=I$ and the impulse responses $\Psi_k$;
      which shock has the longest-lasting effect on which variable?
\end{enumerate}
\end{exercise}

\begin{exercise}[Dickey--Fuller by simulation (not from the course)]\label{ch09:ex-x-df}
\leavevmode
\begin{enumerate}[(a)]
  \item Simulate $10^4$ random walks of length $500$ and compute $t_{\phi=1}$ from the regression of $X_t$ on $X_{t-1}$ with and without intercept; check the $5\%$ quantiles $-1.95$ and $-2.89$.
  \item Repeat with an AR(1) with $\phi=0.95$ and estimate the power of the test at the $5\%$ level; explain why
      non-rejection on a persistent yield series is weak evidence for a unit root.
  \item Verify the identity $\sum_{t=1}^TX_{t-1}\eta_t=\frac12(X_T^2-\sum\eta_t^2)$ and show that its expectation is zero; then estimate $\E[T(\hat\phi-1)]$ by simulation and explain why the ratio in \eqref{ch09:eq-df} is nevertheless biased downwards (the denominator is positively correlated with the numerator's negative part).
\end{enumerate}
\end{exercise}

\begin{exercise}[Cointegration and the VECM (not from the course)]\label{ch09:ex-x-coint}
For the pair of \cref{ch09:ex-coint}
\begin{enumerate}[(a)]
  \item rewrite the system as a VAR(1) in levels and verify that $I+\Pi$ has eigenvalues $1$ and $0.7$.
  \item Show that $\Pi$ is not invertible and identify $\alpha$ and $\beta$; check that another factorisation $(\alpha G)(\beta G^{-\mathsf T})\T$ gives the same $\Pi$.
  \item Simulate the system, regress $y_t$ on $x_t$ and apply an ADF test to the residuals (Engle--Granger); do the same for two independent random walks and comment on ``spurious regression''.
  \item What changes in the VECM if both $\varepsilon_{1t}$ and $u_{t-1}$ enter the equation of $x_t$?
\end{enumerate}
\end{exercise}

\begin{exercise}[Kalman filter by hand (not from the course)]\label{ch09:ex-x-kf}
Take the local-level model with $Q=1$, $H=4$ and initial state $N(0,10)$.
\begin{enumerate}[(a)]
  \item Run the filter for the observations $y_1=1.5$, $y_2=0.8$, $y_3=2.2$: compute $G_t$, $\mu_{t|t}$, $P_{t|t}$.
  \item Compute the two-step forecast of $y_5$ and its variance.
  \item Show that $P_{t|t-1}$ converges monotonically to
      $(Q+\sqrt{Q^2+4QH})/2$.
  \item Show that $G_t\to0$ as $Q\to0$ and interpret it (the level is constant and the filter reduces to the sample mean).
\end{enumerate}
\end{exercise}

\begin{exercise}[Time-varying beta (not from the course)]\label{ch09:ex-x-beta}
In the state-space model of the time-varying CAPM, write out the Kalman recursions for $\bm s_t=(\alpha_t,\beta_t)\T$.
\begin{enumerate}[(a)]
  \item Show that if $\sigma_\xi=\sigma_\zeta=0$ the recursion reproduces recursive least squares for the simple regression.
  \item Explain how the ratio
      $\sigma_\zeta^2/\sigma_\varepsilon^2$ sets the effective window length of the estimate, and relate it to the local-level example.
  \item With daily data for one stock and the S\&P~500, estimate $\sigma_\zeta^2$ by maximising \eqref{ch09:eq-kfll}.
\end{enumerate}
\end{exercise}
